%% TOSCA_Ueqn_Solvers.tex
%% Comprehensive documentation of the momentum equation solvers
%% implemented in TOSCA (src/ueqn.c, src/include/ueqn.h).
%%
%% Compile with:  pdflatex TOSCA_Ueqn_Solvers.tex  (twice for references)
%%

\documentclass[12pt,a4paper]{article}

% ── Packages ──────────────────────────────────────────────────────────────────
\usepackage[T1]{fontenc}
\usepackage[utf8]{inputenc}
\usepackage{lmodern}
\usepackage[margin=2.5cm]{geometry}
\usepackage{amsmath,amssymb,bm}
\usepackage{mathtools}
\usepackage{booktabs}
\usepackage{array}
\usepackage{xcolor}
\usepackage{listings}
\usepackage{hyperref}
\usepackage{caption}
\usepackage{graphicx}
\usepackage{enumitem}
\usepackage{tcolorbox}
\usepackage{parskip}

% ── Colours ───────────────────────────────────────────────────────────────────
\definecolor{tosca-blue}{RGB}{0,84,166}
\definecolor{tosca-grey}{RGB}{90,90,90}
\definecolor{code-bg}{RGB}{248,248,248}
\definecolor{code-kw}{RGB}{0,0,200}
\definecolor{code-cm}{RGB}{0,130,0}
\definecolor{warn-bg}{RGB}{255,248,220}
\definecolor{warn-border}{RGB}{200,160,0}
\definecolor{info-bg}{RGB}{230,244,255}
\definecolor{info-border}{RGB}{0,100,180}

% ── Listings style ────────────────────────────────────────────────────────────
\lstdefinestyle{cstyle}{
	language=C,
	backgroundcolor=\color{code-bg},
	basicstyle=\ttfamily\small,
	keywordstyle=\color{code-kw}\bfseries,
	commentstyle=\color{code-cm}\itshape,
	stringstyle=\color{red!60!black},
	frame=single,
	framerule=0.4pt,
	rulecolor=\color{tosca-grey!40},
	breaklines=true,
	numbers=left,
	numberstyle=\tiny\color{tosca-grey},
	numbersep=8pt,
	tabsize=4,
	showstringspaces=false,
	morekeywords={PetscErrorCode,PetscReal,PetscInt,PetscBool,Vec,SNES,KSP,PC,
		Mat,DM,DMDALocalInfo,MPI_Comm,word,Cmpnts,
		PETSC_TRUE,PETSC_FALSE,PETSC_NULL,PETSC_DEFAULT}
}
\lstset{style=cstyle}

% ── tcolorbox presets ─────────────────────────────────────────────────────────
\tcbuselibrary{breakable}
\newtcolorbox{notebox}[1][]{
	colback=info-bg, colframe=info-border,
	fonttitle=\bfseries, title=#1,
	breakable, left=6pt, right=6pt}
\newtcolorbox{warnbox}[1][]{
	colback=warn-bg, colframe=warn-border,
	fonttitle=\bfseries, title=#1,
	breakable, left=6pt, right=6pt}

% ── Math macros ───────────────────────────────────────────────────────────────
\newcommand{\bU}{\bm{U}}
\newcommand{\bu}{\bm{u}}
\newcommand{\bn}{\bm{n}}
\newcommand{\bxi}{\bm{\xi}}
\newcommand{\bfp}{\bm{f}_p}
\newcommand{\cN}{\mathcal{N}}
\newcommand{\cL}{\mathcal{L}}
\newcommand{\cP}{\mathcal{P}}
\newcommand{\cS}{\mathcal{S}}
\newcommand{\dt}{\Delta t}
\newcommand{\half}{\tfrac{1}{2}}
\newcommand{\ddt}[1]{\dfrac{\partial #1}{\partial t}}
\newcommand{\ddx}[2]{\dfrac{\partial #1}{\partial #2}}
\newcommand{\doi}[1]{\href{https://doi.org/#1}{\texttt{doi:#1}}}

% ── Hyperref ──────────────────────────────────────────────────────────────────
\hypersetup{
	colorlinks=true,
	linkcolor=tosca-blue,
	urlcolor=tosca-blue,
	citecolor=tosca-blue,
	pdftitle={TOSCA U-equation Solvers},
	pdfauthor={TOSCA development notes}
}

% ──────────────────────────────────────────────────────────────────────────────
\begin{document}
	
	\begin{titlepage}
		\centering
		\vspace*{2cm}
		{\LARGE\bfseries\color{tosca-blue} TOSCA Momentum Equation ---\\[0.4em]
			Solvers, Discretisation, and Implementation}\\[1.5em]
		{\large Development notes for \texttt{src/ueqn.c} and \texttt{src/include/ueqn.h}}\\[3em]
		\hrule height 1pt
		\vspace{1em}
		{\large March 2026}\\[0.5em]
		{\normalsize Covers all \texttt{ddtScheme} options:\\
			\texttt{crankNicholson} $\cdot$ \texttt{forwardEuler} $\cdot$
			\texttt{rungeKutta4} $\cdot$ \texttt{IMEX}}
		\vspace{1em}
		\hrule height 1pt
		\vfill
		{\small This document serves as a self-contained reference for the
			PETSc-based momentum solver in TOSCA, including all recent
			optimisations (loop fusion, KSP selector, IMEX-CNAB scheme).}
	\end{titlepage}
	
	\tableofcontents
	\clearpage
	
	% ==============================================================================
	\section{Overview of TOSCA}
	% ==============================================================================
	
	TOSCA is a curvilinear, incompressible Navier--Stokes solver built on PETSc.
	It uses a \emph{staggered} arrangement of contravariant fluxes on a
	body-fitted structured mesh. The governing equations are solved on a
	computational domain $\hat\Omega$ obtained by the diffeomorphism
	$\bm{x}(i,j,k)$, where $(i,j,k)$ are the curvilinear indices.
	
	The primary velocity unknown is the \emph{contravariant flux}
	\begin{equation}
		U_\alpha = J^{-1}\, g^{\alpha\beta}\, u_\beta / J
		\qquad \alpha = 1,2,3
	\end{equation}
	stored in \texttt{ueqn->Ucont}.
	The auxiliary Cartesian velocity \texttt{ueqn->Ucat}
	is reconstructed via \texttt{contravariantToCartesian} and
	is used for viscous stress and Coriolis evaluations.
	
	% ==============================================================================
	\section{Governing Equations}
	% ==============================================================================
	
	\subsection{Incompressible Navier--Stokes in Contravariant Form}
	
	The incompressible momentum equation in strong conservation form reads
	\begin{equation}
		\ddt{\bu} + \nabla\cdot(\bu\otimes\bu)
		= -\nabla p + \nabla\cdot(2\nu_\text{eff}\,\mathbf{S})
		+ \bm{f}
		\label{eq:NS}
	\end{equation}
	where $\nu_\text{eff} = \nu + \nu_t$ is the sum of molecular and SGS
	viscosity (LES), $\mathbf{S}=\tfrac{1}{2}(\nabla\bu+\nabla\bu^\top)$
	is the strain-rate tensor, and $\bm{f}$ collects body forces
	(Coriolis, buoyancy, canopy drag, ABL driving source, wind-farm
	actuator).
	
	After projection onto the curvilinear face-area vectors the momentum
	equation is cast in semi-discrete form as
	\begin{equation}
		\boxed{
			\ddt{U} = \underbrace{-\cP(U)}_{\text{pressure}}
			+ \underbrace{\cN(U)}_{\text{convection}}
			+ \underbrace{\cL(U)}_{\text{viscous}}
			+ \underbrace{\cS}_{\text{sources}}
		}
		\label{eq:semidiscrete}
	\end{equation}
	where each operator is a vector of size $N_\text{dof}=3\,N_x N_y N_z$.
	
	\begin{notebox}[Sign convention]
		The convective operator $\cN(U)$ and viscous operator $\cL(U)$
		are defined to appear on the right-hand side with a positive sign.
		The divergence contribution in \texttt{FormU} is stored as
		$-\nabla\cdot(\bu\,U)$, so a positive contribution to the RHS
		corresponds to acceleration.
	\end{notebox}
	
	\subsection{Source Terms}
	
	\begin{center}
		\begin{tabular}{lll}
			\toprule
			Term & Symbol & Activation flag \\
			\midrule
			Pressure gradient & $\cP = \nabla p$ & always active \\
			Coriolis & $\bm{f}_\text{c} = -2\bm{\Omega}\times\bu$ & \texttt{isAblActive}, \texttt{coriolisActive} \\
			Buoyancy & $\bm{f}_\theta = g\,\theta'/\theta_0\,\hat{\bm{e}}_z$ & \texttt{isTeqnActive} \\
			Canopy drag & $\bm{f}_y = -C_{ft}|\bu|^2/(2H_c)\,\hat{\bm{d}}$ & \texttt{isCanopyActive} \\
			ABL driving & $\bm{f}_s$ (PI controller) & \texttt{controllerActive} \\
			Wind-farm & \texttt{sourceFarmCont} & \texttt{isWindFarmActive} \\
			Sponge/Rayleigh & \texttt{dampingSourceU} & various damping flags \\
			\bottomrule
		\end{tabular}
	\end{center}
	
	% ==============================================================================
	\section{\texttt{FormU}: the Core Assembly Routine}
	% ==============================================================================
	
	\texttt{FormU} is the workhorse that evaluates the convective and
	viscous contributions to the momentum RHS. Its signature is:
	
	\begin{lstlisting}
		PetscErrorCode FormU(
		ueqn_    *ueqn,
		Vec      &Rhs,            // output RHS vector (accumulated)
		PetscReal scale,          // weight applied to all contributions
		PetscInt  formMode = 0);  // 0=full, 1=conv-only, 2=visc-only
	\end{lstlisting}
	
	\subsection{Internal Algorithm}
	
	\textbf{Step 1 -- Divergence and viscous local arrays.}
	The function zeroes the local arrays
	\texttt{lDiv1/2/3} and \texttt{lVisc1/2/3} and populates them in a
	single fused loop over all interior cell faces
	(i, j, k directions computed in one pass):
	\begin{itemize}
		\item \emph{Divergence} (\texttt{formMode $\neq$ 2}):
		upwind or central interpolation of Cartesian velocity
		$\bu$ to each face, multiplied by the face-normal contravariant
		flux (flux limiter optionally applied).
		Available schemes: \texttt{centralUpwind}, \texttt{centralUpwindW},
		\texttt{central}, \texttt{central4}, \texttt{weno3},
		\texttt{quickDiv}.
		\item \emph{Viscous stress} (\texttt{formMode $\neq$ 1}):
		velocity gradients computed via curvilinear chain rule and
		projected onto the face metric tensor $g^{\alpha\beta}$;
		includes the SGS contribution $\nu_t$ when LES is active.
	\end{itemize}
	
	\textbf{Step 2 -- Halo exchange.}
	\texttt{DMLocalToLocal} scatters the populated local arrays to ghost
	cells so that the subsequent divergence stencil can access them.
	
	\textbf{Step 3 -- Cell-centred flux assembly.}
	A second loop assembles the face contributions into the cell-centred
	field \texttt{lFp}:
	\begin{equation}
		f_{p,i} = \sum_{\alpha=1}^{3}
		\Bigl[\texttt{div}_\alpha[k][j][i] - \texttt{div}_\alpha[k-\delta_z][j-\delta_y][i-\delta_x]\Bigr]
		+ [\text{viscous analogously}]
	\end{equation}
	Boundary ghost-cell corrections (no-slip, symmetry) are applied here.
	
	\textbf{Step 4 -- Projection loop.}
	Each cell-centred flux $f_p$ is projected onto the three face-area
	vectors to obtain the contravariant contribution to \texttt{Rhs}:
	\begin{align}
		\texttt{rhs}[k][j][i].x &\mathrel{+}=
		\texttt{scale} \cdot \texttt{iaj}[k][j][i]
		\cdot \frac{1}{2}\bigl[
		(\bm{\xi}_{i}\cdot \bm{f}_p)[k][j][i]
		+ (\bm{\xi}_{i}\cdot \bm{f}_p)[k][j][i{+}1]
		\bigr] \\
		\texttt{rhs}[k][j][i].y &\mathrel{+}=
		\texttt{scale} \cdot \texttt{jaj}[k][j][i]
		\cdot \frac{1}{2}\bigl[
		(\bm{\eta}_{j}\cdot \bm{f}_p)[k][j][i]
		+ (\bm{\eta}_{j}\cdot \bm{f}_p)[k][j{+}1][i]
		\bigr] \\
		\texttt{rhs}[k][j][i].z &\mathrel{+}=
		\texttt{scale} \cdot \texttt{kaj}[k][j][i]
		\cdot \frac{1}{2}\bigl[
		(\bm{\zeta}_{k}\cdot \bm{f}_p)[k][j][i]
		+ (\bm{\zeta}_{k}\cdot \bm{f}_p)[k{+}1][j][i]
		\bigr]
	\end{align}
	
	\subsection{\texttt{formMode} Parameter}
	
	\begin{center}
		\begin{tabular}{cll}
			\toprule
			\texttt{formMode} & Included physics & Use case \\
			\midrule
			0 (default) & convection + viscosity & all legacy callers, explicit schemes \\
			1 & convection only & IMEX: compute $\cN^n$ explicitly \\
			2 & viscosity only  & IMEX: evaluate $\cL(U^{n+1})$ implicitly \\
			\bottomrule
		\end{tabular}
	\end{center}
	
	\begin{notebox}[Default argument]
		The C++ default argument \texttt{formMode = 0} ensures that all
		existing calls (\texttt{FormExplicitRhsU}, \texttt{main.c},
		\texttt{precursor.c}) continue to work without modification.
	\end{notebox}
	
	% ==============================================================================
	\section{Time-Stepping Schemes}
	% ==============================================================================
	
	TOSCA supports four time-stepping schemes selected via
	\texttt{-dUdtScheme} in \texttt{control.dat}. All schemes operate on
	the contravariant flux vector $U = $\texttt{Ucont}.
	
	% ------------------------------------------------------------------------------
	\subsection{Forward Euler (\texttt{forwardEuler})}
	% ------------------------------------------------------------------------------
	
	\paragraph{Method.}
	First-order explicit (Euler forward):
	\begin{equation}
		U^{n+1} = U^n + \dt\,\mathcal{F}(U^n)
	\end{equation}
	where $\mathcal{F} = \cN + \cL + \cS - \cP$ is the full RHS.
	
	\paragraph{Implementation (\texttt{UeqnEuler}).}
	\begin{enumerate}
		\item \texttt{FormExplicitRhsU} evaluates $\mathcal{F}(U^n)$ into \texttt{Rhs}.
		\item \texttt{VecScale(Rhs, dt)}, then \texttt{VecAXPY(Rhs, 1, Ucont)}.
		\item Solution stored: \texttt{VecCopy(Rhs, Ucont)}.
	\end{enumerate}
	
	\paragraph{Stability.}
	Conditionally stable. Requires $\text{CFL} < 1$ and
	$\nu\,\dt/\Delta x^2 < \tfrac{1}{2}$ (viscous CFL). Practical for
	inviscid problems at moderate CFL.
	
	% ------------------------------------------------------------------------------
	\subsection{Runge--Kutta 4 (\texttt{rungeKutta4})}
	% ------------------------------------------------------------------------------
	
	\paragraph{Method.}
	Classic four-stage, fourth-order explicit RK (RK4):
	\begin{equation}
		U^{n+1} = U^n + \dt\,\sum_{i=1}^{4} b_i K_i,
		\qquad
		K_i = \mathcal{F}(U^n + a_i\,\dt\,K_{i-1})
	\end{equation}
	with Butcher tableau coefficients
	$\bm{b} = (\tfrac{1}{6},\tfrac{1}{3},\tfrac{1}{3},\tfrac{1}{6})^\top$
	and $\bm{a} = (0,\tfrac{1}{2},\tfrac{1}{2},1)^\top$.
	
	\paragraph{Implementation (\texttt{UeqnRK4}).}
	Four sequential calls to \texttt{FormExplicitRhsU} (each evaluates
	$\mathcal{F}$ via \texttt{FormU(..., PETSC\_FALSE, 0)}). Accumulation
	into \texttt{Utmp}; result copied to \texttt{Ucont}.
	
	\paragraph{Cost.} Four \texttt{FormU} evaluations per step
	($\approx 2\times$ the cost of forward Euler per equivalent accuracy).
	Best choice when the time step is limited by accuracy rather than
	stability.
	
	% ------------------------------------------------------------------------------
	\subsection{Crank--Nicolson / SNES (\texttt{crankNicholson})}
	\label{sec:cranknicholson}
	% ------------------------------------------------------------------------------
	
	\paragraph{Method.}
	Fully implicit, first-order accurate (second-order in practice with the
	Adams--Bashforth blending of $\texttt{Rhs\_o}$):
	\begin{equation}
		\frac{U^{n+1} - U^n}{\dt}
		= \mathcal{F}(U^{n+1})
		+ \half\, \mathcal{F}(U^n)
		\qquad [\text{it} > \text{itStart}]
		\label{eq:BE}
	\end{equation}
	This is a nonlinear equation for $U^{n+1}$; it is solved using
	PETSc's SNES framework.
	
	\subsubsection{SNES Configuration}
	
	\begin{center}
		\begin{tabular}{ll}
			\toprule
			Component & Setting \\
			\midrule
			SNES type & \texttt{SNESNEWTONTR} (Newton trust-region) \\
			Jacobian  & \texttt{MATMFFD\_WP} (matrix-free finite difference) \\
			Eisenstat--Walker & EW3 (adaptive inner KSP tolerance) \\
			Max outer iterations & \texttt{-snesMaxItersU} (default: 20) \\
			Abs.\ exit tolerance & \texttt{-absTolU} (default: $10^{-5}$) \\
			Rel.\ exit tolerance & \texttt{-relTolU} (default: $10^{-30}$) \\
			\midrule
			KSP solver & \texttt{-kspTypeU}: \texttt{BiCGStab} (default) or \texttt{GMRES} \\
			GMRES restart & \texttt{-kspGMRESRestartU} (default: 30) \\
			Preconditioner & \texttt{PCNONE} \\
			\bottomrule
		\end{tabular}
	\end{center}
	
	\subsubsection{SNES Residual: \texttt{UeqnSNES}}
	
	The residual function registered with PETSc for the Crank--Nicolson
	SNES is:
	\begin{equation}
		F(U) = \dt\,\Bigl[\bm{b}_U + \texttt{FormU}(U,\,\texttt{scale})\Bigr]
		- U + U^n = 0
		\label{eq:SNESres}
	\end{equation}
	where $\bm{b}_U$ is a precomputed constant-RHS buffer (built once per
	time step in \texttt{SolveUEqn}) and \texttt{scale} equals $\half$ for
	$\text{it} > \text{itStart}$ and $1$ otherwise.
	
	\paragraph{Key optimisation: \texttt{bU} buffer.}
	$\bm{b}_U$ accumulates all source terms that do \emph{not} depend on
	the current SNES iterate $U$:
	\begin{equation}
		\bm{b}_U = -\nabla p
		\;\pm\text{buoyancy}
		\;+\; \half\,\mathcal{F}^n
		\;+\;\text{farm source}
	\end{equation}
	This is evaluated \emph{once} before \texttt{SNESSolve} is called,
	rather than on every SNES function evaluation
	($\mathcal{O}(\text{outer}\times\text{inner})$ calls).
	
	\subsubsection{SNES Initial Guess: Explicit Euler Predictor}
	
	Rather than starting from $U^n$, an explicit Euler prediction is used:
	\begin{equation}
		U^{\star} = U^n + \dt\left(\bm{b}_U + \half\,\mathcal{F}^n\right)
		\label{eq:predictor}
	\end{equation}
	implemented as:
	\begin{lstlisting}
		VecCopy  (ueqn->Ucont_o, ueqn->Utmp);
		VecAXPY  (ueqn->Utmp, clock->dt,        ueqn->bU);   // bU holds 0.5*Rhs_o term
		VecAXPY  (ueqn->Utmp, 0.5*clock->dt,    ueqn->Rhs_o);// +other 0.5 -> full Rhs_o weight
		SNESSolve(ueqn->snesU, PETSC_NULL,       ueqn->Utmp);
	\end{lstlisting}
	The predictor typically reduces the number of SNES outer iterations by
	2--3$\times$ compared to starting from $U^n$.
	
	\subsubsection{Procedure Summary (\texttt{SolveUEqn} / Crank--Nicolson)}
	
	\begin{enumerate}
		\item Build $\bm{b}_U$ (pressure, buoyancy, $\half\mathcal{F}^n$, farm).
		\item Compute explicit predictor $U^\star$.
		\item \texttt{SNESSolve} (calls \texttt{UeqnSNES} internally for each residual eval):
		\begin{enumerate}
			\item Copy $\bm{b}_U \to \texttt{Rhs}$.
			\item \texttt{FormU}$(U, \texttt{Rhs}, \text{scale}, 0)$
			(conv + visc + Coriolis + canopy + sourceU).
			\item \texttt{dampingSourceU}, \texttt{meanGradPForcing} (separate traversals).
			\item Scale by $\dt$, add $-U + U^n$.
		\end{enumerate}
		\item Store result: \texttt{VecCopy(Utmp, Ucont)}.
		\item Report: final residual norm, SNES iterations, total linear iterations, wall time.
	\end{enumerate}
	
	% ------------------------------------------------------------------------------
	\subsection{IMEX-CNAB (\texttt{IMEX})}
	\label{sec:imex}
	% ------------------------------------------------------------------------------
	
	\paragraph{Method.}
	Implicit--Explicit Crank--Nicolson Adams--Bashforth (IMEX-CNAB):
	convection is treated \emph{explicitly} with a \emph{blended} Adams--Bashforth
	stencil, while viscosity is treated \emph{implicitly} with
	Crank--Nicolson (CN):
	
	\begin{equation}
		\boxed{
			\frac{U^{n+1} - U^n}{\dt}
			= \bm{b}_U
			+ \underbrace{\half\,\cL(U^{n+1})}_{\text{implicit (CN)}}
		}
		\label{eq:IMEX}
	\end{equation}
	where the explicit buffer is
	\begin{equation}
		\bm{b}_U
		= -\cP
		\;\pm\text{buoy}
		\;+\;\underbrace{\tfrac{3}{2}\cN^n
			- \tfrac{1}{2}\cN^{n-1}}_{\text{AB2 convection}}
		\;+\;\underbrace{\half\cL^n}_{\text{CN explicit half}}
		\;+\;\cS
		\label{eq:bU}
	\end{equation}
	with $\cS$ containing all source terms (Coriolis, canopy, sourceU,
	farm, damping, meanGradP).
	Convection uses Adams--Bashforth~2 (AB2, coefficients $3/2$ and $-1/2$)
	on all steps after the first; on the first step \texttt{RhsConv\_o} is
	not yet available and pure forward-Euler ($\cN^n$ only) is used.

	\begin{notebox}[First step]
		On the first time-step (\texttt{it == itStart}), \texttt{RhsConv\_o} is
		not yet available. The scheme uses pure forward-Euler convection
		($\cN^n$ only) for that step.
	\end{notebox}
	
	\paragraph{Key property: exact linearity and direct KSP.}
	The residual
	\begin{equation}
		F(U^{n+1}) = \dt\,\bm{b}_U
		+ \dt\cdot\text{scale}\cdot\cL(U^{n+1})
		- U^{n+1} + U^n = 0
		\label{eq:IMEXres}
	\end{equation}
	is \emph{linear} in $U^{n+1}$ because $\cL$ (the discrete viscous operator)
	is a linear function of its argument.
	Rearranging gives the linear system
	\begin{equation}
		A\,U^{n+1} = b,
		\qquad
		A\,v \;\equiv\; v - \dt\cdot\text{scale}\cdot\cL(v),
		\qquad
		b \;=\; U^n + \dt\,\bm{b}_U,
		\label{eq:IMEXlinear}
	\end{equation}
	where $\text{scale} = \half$ for it $>$ itStart and $1$ on the first step.
	TOSCA exploits this linearity by solving~\eqref{eq:IMEXlinear}
	\emph{directly} with a standalone PETSc \texttt{KSP} object
	(\texttt{kspIMEX}), bypassing SNES entirely and eliminating the Newton
	trust-region infrastructure.

	\paragraph{Stability advantage.}
	The CN viscosity is unconditionally stable, eliminating the viscous
	time-step restriction $\nu\,\dt/\Delta x^2 < \tfrac{1}{2}$ that plagues
	fully explicit schemes at high Reynolds number or fine mesh spacing.
	The maximum admissible CFL depends on $\beta$ (see the table above):
	pure AB2 ($\beta=1$) is limited to CFL $\lesssim 0.5$, while forward
	Euler ($\beta=0$) tolerates CFL $<1$ at the cost of first-order temporal
	accuracy in the convective term.
	
	\paragraph{Accuracy.}
	With $\beta=1$ (AB2 + CN) the scheme is formally second-order in time.
	With $\beta<1$ the convective part is a blended first/second-order method;
	the overall scheme degrades to first order for $\beta=0$.
	On the first step the scheme is always first-order regardless of $\beta$
	(see note above).
	
	\subsubsection{The \texttt{IMEXMatVec} MatShell Operator}
	\label{sec:imexmatvec}

	The linear operator $A$ in~\eqref{eq:IMEXlinear} is registered with PETSc
	as a \texttt{MatShell} (matrix-free) object stored in \texttt{ueqn->JvIMEX}.
	Each matrix-vector product $w = A\,v$ is performed by the callback
	\texttt{IMEXMatVec}:

	\begin{lstlisting}
PetscErrorCode IMEXMatVec(Mat A, Vec v, Vec Av)
{
    PetscReal dt    = clock->dt;
    PetscReal scale = (clock->it > clock->itStart) ? 0.5 : 1.0;

    // 1. Load v into the Ucont workspace and scatter to ghost cells
    VecCopy(v, ueqn->Ucont);
    DMGlobalToLocalBegin/End(fda, Ucont -> lUcont);

    // 2. Apply boundary conditions to the iterate v
    resetNoPenetrationFluxes(ueqn);      // zero no-slip/slip walls
    resetFacePeriodicFluxesVector(...);

    // 3. Reconstruct Cartesian velocity for viscous stress evaluation
    contravariantToCartesian(ueqn);
    resetCellPeriodicFluxes(...);

    // 4. Evaluate scale * L(v) into scratch buffer ueqn->Rhs
    VecSet(ueqn->Rhs, 0.0);
    FormU(ueqn, ueqn->Rhs, scale, 2); // visc-only: Rhs = scale*L(v)
    VecScale(ueqn->Rhs, dt);                        // Rhs = dt*scale*L(v)
    resetNonResolvedCellFaces(mesh, ueqn->Rhs);     // zero non-resolved faces in operator

    // 5. Av = v - dt*scale*L(v)   (identity minus viscous correction)
    VecCopy(v, Av);
    VecAXPY(Av, -1.0, ueqn->Rhs);
}
	\end{lstlisting}

	\begin{notebox}[Scratch vectors and aliasing]
		\texttt{IMEXMatVec} uses \texttt{ueqn->Ucont}, \texttt{ueqn->lUcont},
		\texttt{ueqn->Ucat}, \texttt{ueqn->lUcat}, and \texttt{ueqn->Rhs} as
		workspace on every mat-vec call. The KSP solution is therefore held in
		\texttt{RhsConv\_o} (not in \texttt{Ucont}) to prevent aliasing between
		the KSP iterate and the operator scratch space. After \texttt{KSPSolve}
		completes, \texttt{RhsConv\_o} is copied to \texttt{Ucont} (Step~6,
		Section~\ref{sec:imexproc}).
	\end{notebox}

	\subsubsection{Boundary-Face Sanitization of the KSP Right-Hand Side}
	\label{sec:imexrhs}

	At non-resolved faces (unused ghost faces, IBM solid-fluid interfaces,
	and inflow/outflow boundary faces) the physical flux is either zero or
	prescribed externally and must not be perturbed by source terms.
	In the old \texttt{UeqnIMEXSNES} residual,
	\texttt{resetNonResolvedCellFaces} acted on the combined
	$\dt\,(\bm{b}_U + \text{scale}\,\cL(U^{n+1}))$ vector, which implicitly
	zeroed the $\dt\,\bm{b}_U$ contribution at those faces.
	In the direct-KSP formulation the RHS $b = U^n + \dt\,\bm{b}_U$
	must be sanitized \emph{before} \texttt{KSPSolve} is called:

	\begin{lstlisting}
// b = U^n + dt*bU
VecCopy(ueqn->Ucont_o, ueqn->Utmp);
VecAXPY(ueqn->Utmp, clock->dt, ueqn->bU);
// Zero dt*bU contributions at non-resolved boundary / IBM faces
resetNonResolvedCellFaces(mesh, ueqn->Utmp);
	\end{lstlisting}

	\begin{warnbox}[Blow-up without RHS sanitization]
		Without this call, pressure-gradient and source-term contributions
		accumulate at ghost/IBM/boundary faces each step at rate
		$O(\dt\,|\bm{b}_U|)$, causing exponential growth of the velocity
		field within $O(1/\dt)$ steps.
	\end{warnbox}

	\subsubsection{Reference Residual Formulation: \texttt{UeqnIMEXSNES}}

	The function \texttt{UeqnIMEXSNES} expresses the same linear system
	as a SNES residual and is retained for testing and verification.
	It can be used to cross-check \texttt{IMEXMatVec} by confirming
	$F(U) = b - A\,U$:

	\begin{lstlisting}
		// Copy prebuilt explicit buffer
		VecCopy(ueqn->bU, Rhs);

		// Add ONLY the implicit viscous half (visc-only FormU)
		PetscReal scale = (clock->it > clock->itStart) ? 0.5 : 1.0;
		FormU(ueqn, Rhs, scale, 2);  // formMode=2: viscosity only

		// Form residual F = dt * Rhs - U^{n+1} + U^n
		VecScale(Rhs, dt);
		VecAXPY (Rhs, -1.0, Ucont);
		VecAXPY (Rhs,  1.0, ueqn->Ucont_o);
	\end{lstlisting}

	\subsubsection{Procedure Summary (\texttt{UeqnIMEX})}
	\label{sec:imexproc}
	
	\begin{enumerate}
		\item Synchronise $U^n$ state (ghost cells, periodic BCs, Cartesian).
		\item Evaluate convective RHS at $U^n$:
		$\cN^n \to \texttt{RhsConv}$
		via \texttt{FormU}$(\cdot, 1.0, 1)$ (\texttt{formMode}=1).
		\item Evaluate viscous RHS at $U^n$:
		$\cL^n \to \texttt{Rhs}$
		via \texttt{FormU}$(\cdot, 1.0, 2)$ (\texttt{formMode}=2).
		\item Build $\bm{b}_U$ (see \eqref{eq:IMEX}):
		\begin{lstlisting}
			VecCopy (ueqn->dP, ueqn->bU);    VecScale(ueqn->bU, -1.0);
			// buoyancy (if it > itStart)
			// AB2: VecAXPY(bU, +1.5, RhsConv); VecAXPY(bU, -0.5, RhsConv_o);
			// CN explicit half: VecAXPY(bU, +0.5, Rhs);   (Rhs holds Visc^n)
			// Coriolis, CanopyForce, sourceU, dampingSourceU, meanGradPForcing
		\end{lstlisting}
		\item Build the KSP right-hand side and set the initial guess:
		\begin{lstlisting}
// b = U^n + dt*bU  (stored in Utmp)
VecCopy (ueqn->Ucont_o, ueqn->Utmp);
VecAXPY (ueqn->Utmp, clock->dt, ueqn->bU);
// Zero dt*bU at non-resolved boundary / IBM faces (Section~\ref{sec:imexrhs})
resetNonResolvedCellFaces(mesh, ueqn->Utmp);
// x0 = b  (explicit-Euler is a near-exact initial guess for nearly-identity A)
VecCopy (ueqn->Utmp, ueqn->RhsConv_o);
		\end{lstlisting}
		\item \texttt{KSPSolve(kspIMEX, Utmp, RhsConv\_o)} --- each
		matrix-vector product calls \texttt{IMEXMatVec}
		(Section~\ref{sec:imexmatvec}).
		The solution is returned in \texttt{RhsConv\_o} to avoid aliasing
		with the operator scratch space in \texttt{Ucont}.
		\item \texttt{VecCopy(RhsConv\_o, Ucont)}; scatter to ghost cells.
		\item Rotate convective buffers: \texttt{VecCopy(RhsConv, RhsConv\_o)}.
		\item Report residual norm, KSP iteration count, wall time.
	\end{enumerate}

	\subsubsection{Cost Comparison}
	
	\begin{center}
		\begin{tabular}{lccl}
			\toprule
			Scheme & \texttt{FormU} calls/step & Linear solves & Note \\
			\midrule
			\texttt{forwardEuler}  & 1 & 0 & 1st order, CFL$<$1 \\
			\texttt{rungeKutta4}   & 4 & 0 & 4th order, CFL$<$1 \\
			\texttt{crankNicholson} & $\sim$55 & $\sim$55 & 2nd eff., unconditional \\
			\texttt{IMEX}          & 3 & 1 & 2nd order, CN visc stable \\
			\bottomrule
		\end{tabular}
	\end{center}
	
	The \texttt{crankNicholson} estimate assumes 5 SNES outer iterations
	$\times$ 11 KSP inner iterations; each KSP iteration calls
	\texttt{FormU} twice (via \texttt{MATMFFD} Jacobian-vector product).

	\subsubsection{KSP Solver Choice for IMEX: GMRES vs.\ BiCGStab}
	\label{sec:ksp-theory}

	\paragraph{Spectrum of the IMEX operator.}
	The operator $A\,v = v - \dt\cdot\text{scale}\cdot\cL(v)$, where
	$\cL = \nu_\text{eff}\nabla^2$ (plus SGS contributions), has eigenvalues
	\begin{equation}
		\lambda_A = 1 + \dt\cdot\text{scale}\cdot\nu_\text{eff}\,k^2 \;\geq\; 1,
	\end{equation}
	since the discrete viscous operator $\cL$ is negative semi-definite on
	bounded domains (eigenvalues $\lambda_\cL = -\nu_\text{eff}k^2 \leq 0$).
	$A$ is therefore \emph{positive definite} with all eigenvalues $\geq 1$.
	The condition number
	\begin{equation}
		\kappa(A) \;\approx\;
		\frac{1 + \dt\cdot\text{scale}\cdot\nu_\text{eff}\,k_{\max}^2}{1}
	\end{equation}
	can reach $O(10^2)$--$O(10^3)$ in ABL simulations with large $\dt$ and
	large LES eddy viscosity $\nu_\text{eff} \sim O(1)\;\text{m}^2/\text{s}$.
	Despite positive definiteness, the discrete $A$ exhibits mild
	\emph{non-symmetry} because the contravariant--Cartesian--contravariant
	projection on non-uniform curvilinear meshes is not self-adjoint.

	\paragraph{GMRES.}
	GMRES~\cite{Saad1986} minimises the residual over the Krylov space of dimension $k$:
	$\|r_k\|_2 = \min_{p \in \mathcal{P}_k,\,p(0)=1}\|p(A)\,r_0\|_2$.
	The residual decreases \emph{monotonically}; the method is
	unconditionally convergent for any nonsingular $A$ (within finite
	arithmetic precision). With the default restart $m=30$, GMRES
	handles the mild non-symmetry and moderate conditioning of the IMEX
	operator reliably. \textbf{GMRES is the recommended choice for IMEX.}

	\paragraph{Standard BiCGStab and the breakdown problem.}
	BiCGStab (\texttt{KSPBCGS}) uses Petrov-Galerkin bi-orthogonalisation
	with a shadow residual $\hat{r} = r_0$ (fixed).
	Each step requires the inner product $\langle\hat{r},\,r_k\rangle \neq 0$.
	For the IMEX operator:
	\begin{itemize}
		\item The initial guess $x_0 = b$ is already close to the solution
		      (the operator is nearly the identity for small $\dt\nu_\text{eff}k^2$).
		\item The initial residual $r_0 = b - A\,x_0$ is therefore $O(\dt)$ and
		      may be small relative to machine precision at fine grid spacings.
		\item On non-uniform curvilinear meshes the operator has mild complex eigenvalues,
		      which can cause $\langle\hat{r},\,r_k\rangle \to 0$ (near-breakdown)
		      after a few steps, stalling convergence.
	\end{itemize}
	In practice, \texttt{KSPBCGS} stalls after $O(10)$--$O(100)$ iterations
	without converging, even when the system is well-posed.

	\paragraph{BiCGStab($L$) and the TOSCA implementation.}
	BiCGStab($L$), introduced by Sleijpen and Fokkema~\cite{Sleijpen1993}, appends a polynomial
	smoother of degree $L$ in $A$ to each outer BiCGStab step.
	This yields $2L$ matrix-vector products per outer iteration (same as
	$L$ standard BiCGStab steps) but with superior orthogonality, eliminating
	the near-breakdown modes.
	With $L=2$, the method is robust for the IMEX operator spectrum
	without incurring the large memory footprint of long-restart GMRES
	($m \gg 30$).
	TOSCA automatically promotes \texttt{-kspTypeU BiCGStab} to
	\texttt{KSPBCGSL} with $L=2$ when the \texttt{IMEX} scheme is active:
	\begin{lstlisting}
if(kspType == "BiCGStab" && dUdtScheme == IMEX)
{
    KSPSetType(kspIMEX, KSPBCGSL);
    KSPBCGSLSetEll(kspIMEX, 2);   // L=2 polynomial smoother
}
	\end{lstlisting}

	\paragraph{Tolerance configuration.}
	The IMEX KSP is configured with
	\begin{lstlisting}
KSPSetTolerances(kspIMEX,
    ueqn->relExitTol,   // rtol = 1e-30 (effectively disabled)
    ueqn->absExitTol,   // atol = 1e-5  (governs convergence)
    1e30,               // dtol: no divergence check
    1000);              // max iterations
	\end{lstlisting}
	Setting \texttt{relExitTol} $= 10^{-30}$ means the solver declares
	convergence solely on \texttt{absExitTol}, regardless of the initial
	residual magnitude.
	This is identical to the \texttt{crankNicholson} inner-KSP behaviour and
	avoids excessive iteration when the initial residual is already near the
	absolute tolerance.

	% ==============================================================================
	\section{IMEX Biharmonic Hyperviscosity}
	\label{sec:hypervisc}
	% ==============================================================================

	\subsection{Motivation}

	Central-difference convection schemes (\texttt{central}, \texttt{centralUpwind})
	do not provide sufficient odd-even damping to suppress high-wavenumber
	errrors that accumulate over many time-steps on structured meshes.
	The worst-offender wavenumber is the mesh Nyquist mode
	$\kappa h = \pi$ (i.e.\ a chequerboard pattern with alternating
	signs on adjacent cells). Because this mode is invisible to
	centred second-order operators (zero convective flux), it grows
	unbounded without dedicated dissipation. The symptoms in practice
	are spatial noise, ringing near stiff gradients, and, in
	severe cases, divergence of the SNES solve.

	The standard remedy~\cite{Lele1992,Ghosal1996} for finite-difference
	and finite-volume LES is an \emph{explicit biharmonic}
	(fourth-order Laplacian) correction: it targets energy only near
	$\kappa h = \pi$ while leaving the resolved energy-containing
	motions virtually unaffected.

	\subsection{Mathematical Formulation}

	\subsubsection{Index-Space Biharmonic Operator}

	Define the one-dimensional fourth-order difference operator acting
	on a grid function $v_n$ with uniform index spacing:
	\begin{equation}
		\delta^4_d\, v_n
		= v_{n+2} - 4v_{n+1} + 6v_n - 4v_{n-1} + v_{n-2}.
		\label{eq:delta4}
	\end{equation}
	Note that $\delta^4_d = (\delta^2_d)^2$ where
	$\delta^2_d v_n = v_{n+1} - 2v_n + v_{n-1}$ is the standard
	finite-difference Laplacian.

	The Fourier symbol of $\delta^4_d$ on a grid with spacing $h$ is
	\begin{equation}
		\widehat{\delta^4}(\kappa)
		= 4\bigl(\cos\kappa h - 1\bigr)^2
		\in [0,\,16].
		\label{eq:symbol}
	\end{equation}
	The symbol vanishes at $\kappa h = 0$ (zero effect on the mean)
	and reaches its maximum value of $16$ at the Nyquist wavenumber
	$\kappa h = \pi$. The three-directional sum has the symbol
	$\widehat{\delta^4}_i + \widehat{\delta^4}_j + \widehat{\delta^4}_k$
	with maximum $48$ at the grid corner $(\pi,\pi,\pi)$.

	\subsubsection{Per-Step Velocity Correction}

	IMEX builds an explicit RHS buffer $\bm{b}_U$ (see Section~\ref{sec:imex})
	and then scales it by $\dt$. The hyperviscosity correction is added to this
	buffer as
	\begin{equation}
		\bm{b}_U \mathrel{-}=
		\frac{\varepsilon_4}{\dt}
		\bigl(\delta^4_i + \delta^4_j + \delta^4_k\bigr)\, U^n,
		\label{eq:bU_hv}
	\end{equation}
	so that after the \texttt{VecScale(Rhs, dt)} step the net
	per-step velocity correction is
	\begin{equation}
		\boxed{
			\Delta U_{\mathrm{hv}}
			= -\varepsilon_4
			\bigl(\delta^4_i + \delta^4_j + \delta^4_k\bigr)\, U^n.
		}
		\label{eq:DU_hv}
	\end{equation}
	The coefficient $\varepsilon_4 \geq 0$ is the hyperviscosity amplitude
	set via \texttt{-imexHyperVisc4U} in \texttt{control.dat}.

	\subsubsection{Filtering Interpretation}

	Equation~\eqref{eq:DU_hv} is equivalent to applying the spectral-space
	\emph{low-pass filter}
	\begin{equation}
		\widehat{G}(\bm{\kappa})
		= 1 - \varepsilon_4
		\sum_{d\in\{i,j,k\}} 4(\cos\kappa_d h - 1)^2
		\label{eq:filter}
	\end{equation}
	to the solution at every time step. For small wavenumbers
	$\kappa h \ll 1$, $\widehat{G} \approx 1 - O((\kappa h)^4)$, so
	resolved-scale turbulence is unaffected. At the Nyquist point
	$\bm{\kappa}h = (\pi,\pi,\pi)$, $\widehat{G} = 1 - 48\varepsilon_4$.
	This is the same class of spectral-vanishing-viscosity
	approach discussed by Karamanos \& Karniadakis~\cite{Karamanos2000}
	and the compact-filter framework of Lele~\cite{Lele1992}, here
	evaluated in index space rather than in physical space.

	\subsection{Stability Analysis}
	\label{sec:hypervisc:stability}

	The hyperviscosity correction is applied \emph{explicitly}.
	Write the Fourier-space update for the worst-case mode:
	\begin{equation}
		\hat{U}^{n+1} = \hat{U}^n\,\widehat{G} = \hat{U}^n
		\bigl(1 - 48\,\varepsilon_4\bigr).
	\end{equation}
	For stable, non-amplifying, monotone damping the amplification factor
	must satisfy $0 \leq 1 - 48\varepsilon_4 \leq 1$, giving
	\begin{equation}
		\boxed{\varepsilon_4 \leq \frac{1}{48} \approx 0.021.}
		\label{eq:stability}
	\end{equation}
	Note that this is an \emph{independent} constraint from the convective
	CFL and the viscous parabolic constraint; it involves only $\varepsilon_4$
	and is independent of $\dt$, mesh spacing, and flow velocity.

	\begin{warnbox}[Stability limit]
		Do not set \texttt{-imexHyperVisc4U} $> 0.021$.
		Values in the range $0.021 < \varepsilon_4 < 0.042$
		will over-damp without causing instability;
		$\varepsilon_4 > 0.042$ amplifies the Nyquist mode and causes
		immediately diverging solutions.
	\end{warnbox}

	\subsection{Boundary Treatment}
	\label{sec:hypervisc:bc}

	The stencil~\eqref{eq:delta4} requires four neighbours ($\pm 1$,
	$\pm 2$). TOSCA has one valid interior ghost cell on each side (at
	$0$ and $mx-1$), but the $\pm 2$ reach requires special treatment
	at cells adjacent to boundaries.

	\subsubsection{\texttt{ii/jj/kk\_periodic} (PETSc native periodic)}

	PETSc fills ghost cells at negative and beyond-maximum indices during
	\texttt{DMGlobalToLocalBegin/End}. All four offsets are used directly.
	No index remapping is needed.

	\subsubsection{\texttt{i/j/k\_periodic} (single-process ghosted periodic)}

	\texttt{resetFacePeriodicFluxesVector} sets
	\texttt{lUcont[0] = lUcont[mx-2]} (left ghost). The right ghost
	\texttt{lUcont[mx-1]} is \emph{not} filled. The valid
	read range is $[0,\,mx-2]$ with period $mx-2$.
	Index remapping near both boundaries is therefore required:

	\begin{center}
		\begin{tabular}{cll}
			\toprule
			cell $i$ & Remapped indices & Reason \\
			\midrule
			$i = 1$ & $i_{-2} \leftarrow mx-3$ & wrap $-1$ mod $(mx-2)$ \\
			$i = mx-2$ & $i_{+2} \leftarrow 2$ & avoid junk ghost $mx-1$ \\
			\bottomrule
		\end{tabular}
	\end{center}

	After the loop, $\texttt{rhs}[k][j][0]$ is overwritten with
	$\texttt{rhs}[k][j][mx-2]$ to propagate the computed
	hyperviscosity correction to the left periodic face.

	\subsubsection{Non-periodic (wall, inlet/outlet)}

	Only one filled ghost cell is available at each non-periodic boundary;
	\texttt{lUcont[mx-1]} contains junk data and must not be read.
	All four offsets are clamped to the valid interior range $[1,\,mx-2]$:
	\begin{equation}
		i_{\pm1} = \mathrm{clamp}(i\pm1,\;1,\;mx-2),
		\qquad
		i_{\pm2} = \mathrm{clamp}(i\pm2,\;1,\;mx-2).
		\label{eq:clamp}
	\end{equation}
	When $i = 1$ the clamp gives $i_{-1} = i_{-2} = 1$; substituting
	into~\eqref{eq:delta4} with the repeated-value padding yields
	\begin{equation}
		\delta^4 v_1 = v_3 - 4v_2 + 6v_1 - 4v_1 + v_1
		= v_3 - 4v_2 + 3v_1,
	\end{equation}
	which evaluates to zero when $v$ is locally constant
	(zero-gradient / slip condition). This is the correct physical
	behaviour at a slip wall: the hyperviscosity correction
	vanishes where the flux is already smooth, so \emph{no artificial
	no-slip boundary condition is introduced}.

	\begin{notebox}[Why not mirror-image padding?]
		A symmetric (mirror) ghost would enforce zero normal gradient
		of the \emph{correction} but would couple the correction to the
		normal velocity component and could create spurious wall-normal
		flux. Nearest-value clamping avoids this and is consistent with
		a slip boundary for all three contravariant flux components.
	\end{notebox}

	\subsection{Implementation Summary}

	The function \texttt{hyperViscosityU(ueqn, Rhs, scale)} performs:

	\begin{enumerate}
		\item Early exit if $\varepsilon_4 = 0$.
		\item Compute effective coefficient
		$\texttt{eps} = \texttt{scale}\cdot \varepsilon_4 / \dt$.
		\item Loop over all interior faces
		$(i,j,k)\in[\texttt{lxs},\texttt{lxe})\times\cdots$:
		\begin{enumerate}
			\item Compute remapped indices \texttt{im2,im1,ip1,ip2}
			(and analogues in $j,k$) according to
			Section~\ref{sec:hypervisc:bc}.
			\item For each face type (\texttt{.x}, \texttt{.y}, \texttt{.z}),
			skip IBM solid cells via \texttt{nvert} check, then accumulate
			\begin{align*}
				b_4 &=
				\underbrace{U_{i+2} - 4U_{i+1} + 6U_i - 4U_{i-1} + U_{i-2}}_{
					\delta^4_i\, U_i}\\
				&\quad
				+\underbrace{U_{j+2} - 4U_{j+1} + 6U_j - 4U_{j-1} + U_{j-2}}_{
					\delta^4_j\, U_j}\\
				&\quad
				+\underbrace{U_{k+2} - 4U_{k+1} + 6U_k - 4U_{k-1} + U_{k-2}}_{
					\delta^4_k\, U_k},
			\end{align*}
			then subtract: $\texttt{rhs}[k][j][i] \mathrel{-}= \texttt{eps}\cdot b_4$.
		\end{enumerate}
	\end{enumerate}

	The call site in \texttt{UeqnIMEX} is:
	\begin{lstlisting}
	if(ueqn->hyperVisc4 > 0.0)
	    hyperViscosityU(ueqn, ueqn->bU, 1.0);
	\end{lstlisting}
	Placed \emph{after} all other explicit contributions are assembled
	into \texttt{bU} and \emph{before} the SNES solve.

	\subsection{Control Parameters}

	\begin{center}
		\begin{tabular}{llll}
			\toprule
			Option & Type & Default & Description \\
			\midrule
			\texttt{-imexHyperVisc4U} & real & 0.0 & Hyperviscosity amplitude $\varepsilon_4$
			(\texttt{IMEX} only) \\
			\bottomrule
		\end{tabular}
	\end{center}

	Example:
	\begin{lstlisting}[language={},style=cstyle]
	-dUdtScheme       IMEX
	-divScheme        centralUpwind
	-kspTypeU         BiCGStab
	-imexHyperVisc4U  0.01     # biharmonic hyperviscosity (recommended: 0.005-0.02)
	\end{lstlisting}

	\subsection{Best Practices}
	\label{sec:hypervisc:bestpractices}

	\begin{enumerate}
		\item \textbf{Start small.} Begin with $\varepsilon_4 = 0.005$ and
		increase only if chequerboard noise remains visible in velocity
		spectra or if the SNES fails to converge. Values
		$\varepsilon_4 \in [0.005,\,0.015]$ are typical for ABL simulations
		with \texttt{centralUpwind}.

		\item \textbf{Respect the stability bound.} Always keep
		$\varepsilon_4 < 1/48 \approx 0.021$ (see~\eqref{eq:stability}).
		Values above this threshold will over-damp resolved scales;
		values above $1/24 \approx 0.042$ cause the Nyquist mode to grow.

		\item \textbf{IMEX only.} The parameter \texttt{-imexHyperVisc4U}
		is read only when \texttt{-dUdtScheme IMEX} is selected.
		It has no effect on \texttt{crankNicholson}, \texttt{forwardEuler},
		or \texttt{rungeKutta4} (the SNES schemes are unconditionally stable
		with respect to odd-even modes via their implicit structure).

		\item \textbf{Not a substitute for a good scheme.} Hyperviscosity
		suppresses Nyquist noise but cannot fix a fundamentally unstable
		convection scheme. If the simulation is underresolved, use a higher
		\texttt{divScheme} (e.g.\ \texttt{weno3}) or coarsen the mesh
		rather than increasing $\varepsilon_4$.

		\item \textbf{Verify with spectra.} After setting $\varepsilon_4$,
		plot the energy spectrum of the streamwise velocity at a representative
		plane. The inertial subrange should follow $E(\kappa) \sim \kappa^{-5/3}$
		without artificial pile-up near $\kappa h = \pi$ and without excessive
		kill-off of the $\kappa^{-5/3}$ cascade.

		\item \textbf{Combination with \texttt{central4}.} The
		\texttt{central4} divergence scheme already applies its own
		hyperviscosity via \texttt{-hyperVisc} (acting on Cartesian velocity
		gradients). Using \texttt{-imexHyperVisc4U} simultaneously is
		redundant and may over-damp. Keep one active at a time.

		\item \textbf{IBM boundaries.} The nvert guard
		(\texttt{nvert[k][j][i] + nvert[k][j][i+1] < 1.0}) skips the
		correction at IBM solid-fluid interfaces. No special action is
		required; the correction automatically vanishes at cells flagged
		as solid.
	\end{enumerate}

	% ==============================================================================
	\section{Key Data Structures}
	% ==============================================================================
	
	\subsection{Primary \texttt{ueqn\_} Vectors}
	
	\begin{center}
		\begin{tabular}{lll}
			\toprule
			Field & Type & Description \\
			\midrule
			\texttt{Ucont} / \texttt{lUcont} & \texttt{Vec} & Global/local contravariant flux $U$\\
			\texttt{Ucont\_o} & \texttt{Vec} & $U^n$ (previous step) \\
			\texttt{Ucat} / \texttt{lUcat} & \texttt{Vec} & Cartesian velocity \\
			\texttt{Utmp} & \texttt{Vec} & Scratch / SNES initial guess \\
			\texttt{Rhs} & \texttt{Vec} & Working RHS accumulator \\
			\texttt{Rhs\_o} & \texttt{Vec} & $\mathcal{F}^n$ (full RHS at prev.\ step, \texttt{crankNicholson}) \\
			\texttt{dP} & \texttt{Vec} & Pressure gradient $\nabla p$ \\
			\texttt{bTheta} & \texttt{Vec} & Buoyancy source $\bm{f}_\theta$ \\
			\texttt{sourceU} / \texttt{lsourceU} & \texttt{Vec} & ABL PI-controller source \\
			\midrule
			\multicolumn{3}{l}{\emph{crankNicholson only}} \\
			\texttt{bU} & \texttt{Vec} & Constant RHS buffer (prebuilt per step) \\
			\midrule
			\multicolumn{3}{l}{\emph{IMEX only}} \\
			\texttt{bU} & \texttt{Vec} & Explicit IMEX buffer (prebuilt per step) \\
			\texttt{RhsConv} & \texttt{Vec} & $\cN^n$ (conv-only RHS, current step) \\
			\texttt{RhsConv\_o} & \texttt{Vec} & $\cN^{n-1}$ / KSP solution buffer (see \S\ref{sec:imexmatvec}) \\
			\texttt{kspIMEX} & \texttt{KSP} & Standalone KSP for the IMEX linear solve \\
			\texttt{JvIMEX} & \texttt{Mat} & MatShell: $A\,v = v - \dt\cdot\text{scale}\cdot\cL(v)$ \\
			\midrule
			\multicolumn{3}{l}{\emph{Intermediate FormU arrays}} \\
			\texttt{lDiv1/2/3} & \texttt{Vec} & Face-local divergence contributions \\
			\texttt{lVisc1/2/3} & \texttt{Vec} & Face-local viscous contributions \\
			\texttt{lFp} & \texttt{Vec} & Cell-centred combined flux (before projection) \\
			\bottomrule
		\end{tabular}
	\end{center}
	
	\subsection{SNES/KSP Objects}

	\paragraph{\texttt{crankNicholson} objects.}
	\begin{center}
		\begin{tabular}{lll}
			\toprule
			Field & Type & Description \\
			\midrule
			\texttt{snesU} & \texttt{SNES} & PETSc nonlinear solver context \\
			\texttt{JU} & \texttt{Mat} & Matrix-free MFFD Jacobian approximation \\
			\texttt{ksp} & \texttt{KSP} & Inner Krylov linear solver (inside SNES) \\
			\texttt{pc} & \texttt{PC} & Preconditioner (currently \texttt{PCNONE}) \\
			\bottomrule
		\end{tabular}
	\end{center}

	\paragraph{\texttt{IMEX} objects.}
	\begin{center}
		\begin{tabular}{lll}
			\toprule
			Field & Type & Description \\
			\midrule
			\texttt{kspIMEX} & \texttt{KSP} & Standalone KSP for the IMEX linear solve \\
			\texttt{JvIMEX} & \texttt{Mat} & MatShell: $A\,v = v - \dt\cdot\text{scale}\cdot\cL(v)$ \\
			\bottomrule
		\end{tabular}
	\end{center}

	\paragraph{Shared solver parameters.}
	\begin{center}
		\begin{tabular}{lll}
			\toprule
			Field & Type & Description \\
			\midrule
			\texttt{kspType} & \texttt{word} & \texttt{"BiCGStab"} or \texttt{"GMRES"} \\
			\texttt{gmresRestart} & \texttt{PetscInt} & GMRES restart parameter (default 30) \\
			\texttt{snesMaxIter} & \texttt{PetscInt} & Max SNES outer iterations (\texttt{crankNicholson}) \\
			\texttt{relExitTol} & \texttt{PetscReal} & Relative stopping tolerance ($10^{-30}$ default) \\
			\texttt{absExitTol} & \texttt{PetscReal} & Absolute stopping tolerance ($10^{-5}$ default) \\
			\bottomrule
		\end{tabular}
	\end{center}
	
	% ==============================================================================
	\section{Control Parameters (\texttt{control.dat})}
	% ==============================================================================
	
	\begin{center}
		\begin{tabular}{llll}
			\toprule
			Key & Type & Default & Description \\
			\midrule
			\texttt{-dUdtScheme} & string & --- & Time scheme (mandatory) \\
			\texttt{-divScheme} & string & --- & Divergence scheme (mandatory) \\
			\texttt{-kspTypeU} & string & --- & KSP type (mandatory if implicit) \\
			\texttt{-snesMaxItersU} & int & 20 & Max SNES outer iterations \\
			\texttt{-absTolU} & real & $10^{-5}$ & SNES absolute tolerance \\
			\texttt{-relTolU} & real & $10^{-30}$ & SNES relative tolerance \\
			\texttt{-kspGMRESRestartU} & int & 30 & GMRES restart (if GMRES) \\
			\texttt{-hyperVisc} & real & 1.0 & Hyperviscosity (if \texttt{central4}) \\
			\texttt{-imexHyperVisc4U} & real & 0.0 & Biharmonic hyperviscosity $\varepsilon_4$ (\texttt{IMEX} only, see \S\ref{sec:hypervisc}) \\
		\end{tabular}
	\end{center}
	
	Example \texttt{control.dat} block:
	
	\begin{lstlisting}[language={},style=cstyle]
		# Momentum equation
		-dUdtScheme     crankNicholson  # or: forwardEuler, rungeKutta4, IMEX
		-divScheme      centralUpwind
		-kspTypeU       BiCGStab        # mandatory for crankNicholson and IMEX
		-snesMaxItersU  5
		-absTolU        1e-5
	\end{lstlisting}
	
	For IMEX (second-order, CFL$\,{<}\,0.5$; GMRES recommended,
	see Section~\ref{sec:ksp-theory}):
	\begin{lstlisting}[language={},style=cstyle]
	-dUdtScheme       IMEX
	-divScheme        centralUpwind
	-kspTypeU         GMRES          # recommended for IMEX
	-kspGMRESRestartU 30
	-absTolU          1e-5
	\end{lstlisting}

	\begin{notebox}[BiCGStab for IMEX]
		Setting \texttt{-kspTypeU BiCGStab} with the IMEX scheme automatically
		activates \texttt{KSPBCGSL} with $L=2$ (BiCGStab with a degree-2
		polynomial smoother) instead of standard \texttt{KSPBCGS}.
		Standard BiCGStab cannot reliably solve the IMEX system ---
		see Section~\ref{sec:ksp-theory} for the breakdown analysis.
	\end{notebox}

	For IMEX at higher CFL (first-order convection, CFL$\,{<}\,1.0$):
	\begin{lstlisting}[language={},style=cstyle]
-dUdtScheme       IMEX
-divScheme        centralUpwind
-kspTypeU         GMRES
	\end{lstlisting}
	
	% ==============================================================================
	\section{Optimisations Implemented}
	% ==============================================================================
	
	\subsection{Precomputed Constant-RHS Buffer (\texttt{bU})}
	
	Prior to this optimisation, every SNES function evaluation recomputed
	all source terms (pressure gradient, buoyancy, wind-farm source) that
	do not depend on the current iterate. For a solve with 5 Newton
	iterations $\times$ 11 KSP steps, this meant $\sim$55 redundant
	evaluations of these terms.
	
	The vector \texttt{bU} is now built \emph{once} per time step in
	\texttt{SolveUEqn} before \texttt{SNESSolve} is called.
	In \texttt{UeqnSNES} the function begins with
	\texttt{VecCopy(ueqn->bU, Rhs)} instead of the sequence
	\texttt{VecSet + multiple VecAXPY}.
	
	\subsection{Loop Fusion in the Projection Step}
	
	The Coriolis, canopy drag, and ABL driving source are evaluated inline
	within the projection loop of \texttt{FormU}, improving data locality
	and eliminating the overhead of three extra loop launches.
	
	\subsection{Dead Code Removal in the Projection Loop}
	
	The original projection loop computed:
	\begin{itemize}
		\item Four dot products per spatial direction ($v0, v1, v2, v3$) to
		support a higher-order stencil.
		\item Twelve cell-width divisions ($d0$--$d3$) for a flux-limited
		scheme.
		\item Three calls to \texttt{getFace2Cell4Stencil*} per cell.
	\end{itemize}
	None of these were used in the final \texttt{central(v1, v2)} assignment.
	Removing them reduces $\sim$40 floating-point operations and 3 function
	calls per cell in the projection loop.
	
	\subsection{Explicit Euler Predictor as SNES Initial Guess}
	
	Starting the Newton iteration from the explicit Euler prediction
	$U^\star$ (\eqref{eq:predictor}) rather than $U^n$ was found to reduce
	the average number of SNES outer iterations from $\sim$8 to $\sim$5
	and total linear iterations from $\sim$25 to $\sim$11 in representative
	ABL test cases.
	
	\subsection{Selectable KSP Solver}

	The Krylov solver is selected at runtime via \texttt{-kspTypeU}.
	The behaviour differs between the two implicit schemes:

	\paragraph{\texttt{crankNicholson} --- SNES inner KSP.}
	Both \texttt{BiCGStab} and \texttt{GMRES} are effective here.
	\texttt{BiCGStab} is the default (lower memory, handles mild non-symmetry).
	\texttt{GMRES} provides monotone residual reduction and is preferred for
	poorly-conditioned or highly non-symmetric cases.
	The Eisenstat--Walker (EW3) adaptive tolerance strategy is always active:
	it progressively tightens the inner KSP tolerance as the outer SNES
	residual decreases, avoiding over-solving in early Newton steps.

	\paragraph{\texttt{IMEX} --- standalone KSP.}
	\textbf{GMRES is strongly recommended} (see Section~\ref{sec:ksp-theory}).
	Standard BiCGStab (\texttt{KSPBCGS}) is automatically replaced by
	\texttt{KSPBCGSL} with $L=2$ to prevent near-breakdown on the IMEX
	near-identity operator.
	The Eisenstadt--Walker strategy does \emph{not} apply here (there is no
	outer SNES).
	Convergence is governed by \texttt{absExitTol}; the relative tolerance
	\texttt{relExitTol} is set to $10^{-30}$ (effectively disabled).
	
	\subsection{Measured Performance (\texttt{crankNicholson})}
	
	\begin{center}
		\begin{tabular}{lccc}
			\toprule
			Version & SNES iters & Linear iters & Wall time / step \\
			\midrule
			Baseline & $\sim$8 & $\sim$25 & 0.717 s \\
			After all optimisations & 5 & 11 & 0.449 s \\
			\midrule
			Improvement & $\mathbf{-37.5\%}$ & -- & $\mathbf{-37\%}$ \\
			\bottomrule
		\end{tabular}
	\end{center}
	
	% ==============================================================================
	\section{Schematics and Flow Diagrams}
	% ==============================================================================
	
	\subsection{\texttt{SolveUEqn} Control Flow}
	
	\begin{center}
		\fbox{
			\begin{minipage}{0.88\textwidth}
				\textbf{SolveUEqn}
				\begin{enumerate}[leftmargin=1.2em]
					\item \texttt{VecSet(Rhs, 0)}
					\item \textbf{if} \texttt{ddtScheme == "crankNicholson"}
					\begin{itemize}
						\item Build \texttt{bU}, compute \texttt{Utmp} (explicit predictor)
						\item \texttt{SNESSolve} $\to$ \texttt{UeqnSNES} (residual) repeated
						\item Print: residual, iters, linear iters, time
					\end{itemize}
					\item \textbf{else if} \texttt{ddtScheme == "forwardEuler"}
					\begin{itemize}
						\item $\to$ \texttt{UeqnEuler}
					\end{itemize}
					\item \textbf{else if} \texttt{ddtScheme == "rungeKutta4"}
					\begin{itemize}
						\item $\to$ \texttt{UeqnRK4}
					\end{itemize}
					\item \textbf{else if} \texttt{ddtScheme == "IMEX"}
					\begin{itemize}
						\item $\to$ \texttt{UeqnIMEX}
					\end{itemize}
					\item Post-processing: \texttt{resetNoPenetrationFluxes},
					IBM BCs, overset BCs, \texttt{adjustFluxesLocal}
				\end{enumerate}
			\end{minipage}
		}
	\end{center}
	
	\subsection{\texttt{FormU} Call Graph}
	
	\begin{center}
		\begin{tabular}{lcc}
			\toprule
			Caller & \texttt{formMode} & Scale \\
			\midrule
			\texttt{UeqnSNES} (crankNicholson) & 0 (full) & $0.5$ or $1$ \\
			\texttt{IMEXMatVec} (IMEX KSP mat-vec) & 2 (visc-only) & $0.5$ or $1$ \\
			\texttt{UeqnIMEXSNES} (reference only) & 2 (visc-only) & $0.5$ or $1$ \\
			\texttt{UeqnIMEX} (conv eval) & 1 (conv-only) & 1 \\
			\texttt{UeqnIMEX} (visc eval) & 2 (visc-only) & 1 \\
			\texttt{FormExplicitRhsU} (fwdEuler/RK4) & 0 (full) & 1 \\
			\texttt{main.c / precursor.c} (init $\mathcal{F}^n$) & 0 (full) & 1 \\
			\bottomrule
		\end{tabular}
	\end{center}
	
	% ==============================================================================
	\section{Known Limitations and Future Work}
	% ==============================================================================
	
	\begin{itemize}
		\item \textbf{Pressure--velocity coupling.} All four schemes use
		\texttt{SolveUEqn} and the pressure solver separately; they do
		not form a coupled problem. The IMEX scheme already uses a direct
		KSP rather than SNES; further cost reduction via a
		fractional-step/projection approach would require a modified
		pressure-correction step to maintain second-order accuracy.
		
		\item \textbf{IMEX first step.} On the first time-step, the IMEX scheme
		uses forward-Euler convection (no $\cN^{n-1}$ available), degrading
		to first-order for that step regardless of $\beta$. A consistent
		initialisation (e.g.\ two smaller sub-steps) would improve transient
		accuracy.
		
		\item \textbf{IMEX CFL constraint.} The pure AB2 ($\beta=1$) stability
		limit CFL $\lesssim 0.5$ can be restrictive. Reducing $\beta$
		relaxes this at the cost of first-order temporal accuracy in the
		convective term. A value $\beta=0$ (forward Euler) is stable for
		CFL $<1$. For CFL $>1$ the fully implicit
		\texttt{crankNicholson} scheme should be used.
		
		\item \textbf{Preconditioner.} \texttt{PCNONE} is used for all
		schemes. A block-ILU(0) or diagonal Jacobi preconditioner could
		reduce the number of KSP iterations from $\sim$11 to $\sim$2--3
		for the Crank--Nicolson scheme.
		
		\item \textbf{Variable time step.} The IMEX bU buffer formula assumes
		constant $\dt$. If variable time stepping is introduced, the AB2
		coefficients must be replaced by BDF2/variable-step AB2.

		\item \textbf{Hyperviscosity anisotropy.} The index-space
		operator~\eqref{eq:delta4} penalises equally in all three index
		directions regardless of the local mesh stretching. On highly
		anisotropic grids (e.g.\ wall-resolved boundary layers with
		$\Delta z / \Delta x \ll 1$), the physical dissipation is
		stronger in the coarsely-resolved wall-normal direction than in
		the streamwise direction. A metric-weighted biharmonic
		($\delta^4 / h^4$) would be more consistent but requires knowledge
		of local mesh spacings and is not currently implemented.
	\end{itemize}
	
	% ==============================================================================
	\appendix
	\section{Relation Between \texttt{Rhs\_o} and \texttt{RhsConv}}
	% ==============================================================================
	
	These two vectors serve different purposes:
	
	\paragraph{\texttt{Rhs\_o} (Crank--Nicolson).}
	Contains the \emph{full} right-hand side $\mathcal{F}(U^n)$
	evaluated by \texttt{FormExplicitRhsU} after each successful step
	in \texttt{main.c}/\texttt{precursor.c}. Used as the Adams--Bashforth
	blend in the Crank--Nicolson bU buffer:
	$\bm{b}_U \mathrel{+}= \half\,\mathcal{F}^n$.
	
	\paragraph{\texttt{RhsConv} / \texttt{RhsConv\_o} (IMEX).}
	Contain \emph{only} the convective contribution $\cN(U^n)$,
	evaluated once per step inside \texttt{UeqnIMEX} before the SNES
	solve. The rotation \texttt{RhsConv} $\to$ \texttt{RhsConv\_o} at
	end of step maintains the AB2 two-step history without external
	intervention from \texttt{main.c}.
	
	\section{Files Modified in This Session}
	
	\begin{center}
		\begin{tabular}{ll}
			\toprule
			File & Changes \\
			\midrule
			\texttt{src/include/ueqn.h} & Added \texttt{RhsConv}, \texttt{RhsConv\_o},
			updated \texttt{bU} doc, \\
			& new \texttt{formMode} default arg, declarations for \\
			& \texttt{UeqnIMEX}, \texttt{UeqnIMEXSNES} \\
			\texttt{src/ueqn.c} & Added \texttt{IMEX} block in \texttt{InitializeUEqn} \\
			& Added \texttt{IMEXMatVec} MatShell callback \\
			& Replaced SNES with standalone \texttt{kspIMEX} + \texttt{JvIMEX} MatShell \\
			& \texttt{resetNonResolvedCellFaces} on KSP RHS (blow-up fix) \\
			& \texttt{KSPBCGSL} ($L$=2) for IMEX; \texttt{rtol}=\texttt{relExitTol} \\
			& Added \texttt{UeqnIMEXSNES} (reference residual formulation) \\
			& Added \texttt{UeqnIMEX} function \\
			& Added \texttt{formMode} param to \texttt{FormU} \\
			& Wired \texttt{UeqnIMEX} into \texttt{SolveUEqn} \\
			& Loop fusion, dead code removal, \texttt{bU} buffer \\
			& Explicit Euler predictor, KSP selector \\
			\texttt{src/main.c} & Updated \texttt{FormU} call signature \\
			\texttt{src/precursor.c} & Updated \texttt{FormU} call signature \\
			\bottomrule
		\end{tabular}
	\end{center}
	
	% ==============================================================================
	\begin{thebibliography}{9}

	\bibitem{Sleijpen1993}
	Sleijpen, G.L.G. and Fokkema, D.R. (1993).
	``BiCGstab($\ell$) for linear equations involving unsymmetric matrices
	with complex spectrum.''
	\textit{Electronic Transactions on Numerical Analysis}, \textbf{1}, 11--32.

	\bibitem{Saad1986}
	Saad, Y. and Schultz, M.H. (1986).
	``GMRES: A generalized minimal residual algorithm for solving nonsymmetric
	linear systems.''
	\textit{SIAM Journal on Scientific and Statistical Computing},
	\textbf{7}(3), 856--869.
	\doi{10.1137/0907058}

	\bibitem{Lele1992}
	Lele, S.K. (1992).
	``Compact finite difference schemes with spectral-like resolution.''
	\textit{Journal of Computational Physics}, \textbf{103}(1), 16--42.
	\doi{10.1016/0021-9991(92)90324-R}

	\bibitem{Ghosal1996}
	Ghosal, S. (1996).
	``An analysis of numerical errors in large-eddy simulations of turbulence.''
	\textit{Journal of Computational Physics}, \textbf{125}(1), 187--206.
	\doi{10.1006/jcph.1996.0088}

	\bibitem{Karamanos2000}
	Karamanos, G.-S. and Karniadakis, G.E. (2000).
	``A spectral vanishing viscosity method for large-eddy simulations.''
	\textit{Journal of Computational Physics}, \textbf{163}(1), 22--50.
	\doi{10.1006/jcph.2000.6552}

	\bibitem{Vreman1996}
	Vreman, B., Geurts, B. and Kuerten, H. (1996).
	``Comparison of numerical schemes in large-eddy simulation of the
	temporal mixing layer.''
	\textit{International Journal for Numerical Methods in Fluids},
	\textbf{22}(4), 297--311.
	\doi{10.1002/fld.1650220403}

	\end{thebibliography}

% ============================================================
%  PART II — TEMPERATURE EQUATION SOLVERS
% ============================================================
\clearpage
\part*{Part II — Temperature Equation Solvers}
\addcontentsline{toc}{part}{Part II — Temperature Equation Solvers}

% ----------------------------------------------------------------
\section{Governing Equation for Temperature}
\label{sec:teqn-governing}
% ----------------------------------------------------------------

The filtered (LES) or Reynolds-averaged temperature transport equation solved by TOSCA is

\begin{equation}
  \frac{\partial T}{\partial t}
  + \nabla\!\cdot\!(\mathbf{u}\,T)
  = \nabla\!\cdot\!(\kappa_{\text{eff}}\,\nabla T)
  + s_T,
  \label{eq:teqn}
\end{equation}

where $T$ is the potential temperature, $\mathbf{u}$ the resolved velocity,
$\kappa_{\text{eff}}$ the effective thermal diffusivity (molecular plus SGS), and
$s_T$ a source/sink term (e.g.\ nudging, Rayleigh damping).
In curvilinear coordinates the convective flux through an $\xi^1$-face reads
$q_{\text{conv}}^{(1)} = -U^1 T$
and the diffusive flux

\begin{equation}
  q_{\text{visc}}^{(1)}
  = \kappa_{\text{eff}}\,\frac{1}{J^2}
    \bigl(g^{11}\,\partial_\xi T
         + g^{21}\,\partial_\eta T
         + g^{31}\,\partial_\zeta T\bigr)
    J,
\end{equation}

with $g^{ij}$ the metric-tensor components and $J$ the Jacobian of the mapping.
The cell-RHS is the sum of all face fluxes times $J$ (stored as \texttt{rhs}).

% ----------------------------------------------------------------
\section{\texttt{FormT} — Core Assembly Routine}
\label{sec:formT}
% ----------------------------------------------------------------

\texttt{FormT(teqn, Rhs, scale, formMode)} assembles the spatial operator
$\mathcal{F}(T) = \nabla\!\cdot\!(-\mathbf{u}T + \kappa_{\text{eff}}\nabla T)$
and writes \texttt{scale}\,$\times \mathcal{F}(T)$ into \texttt{Rhs}.

\subsection*{Signature}

\begin{lstlisting}[language=C,basicstyle=\ttfamily\small]
PetscErrorCode FormT(teqn_ *teqn, Vec Rhs, PetscReal scale,
                     PetscInt formMode);
\end{lstlisting}

\subsection*{\texttt{formMode} parameter}

\begin{center}
\begin{tabular}{clp{8cm}}
\toprule
\textbf{Value} & \textbf{Macro / name} & \textbf{Effect} \\
\midrule
0 & \texttt{FULL} & Both convective and diffusive fluxes assembled. \\
1 & \texttt{CONV\_ONLY} & Only convective fluxes ($\nabla\!\cdot\!(-\mathbf{u}T)$). \\
2 & \texttt{DIFF\_ONLY} & Only diffusive fluxes ($\nabla\!\cdot\!(\kappa_{\text{eff}}\nabla T)$). \\
\bottomrule
\end{tabular}
\end{center}

\subsection*{Algorithm}

\begin{enumerate}
  \item Scatter ghost cells for $T$, $\mathbf{U}^{\text{cont}}$, face normals,
        $\nu_t$ (SGS), limiters.
  \item \textbf{Fused face loop} — a single triple-nested loop over
        $(k,j,i)\in[z_s,z_e)\times[y_s,y_e)\times[x_s,x_e)$:
        \begin{itemize}
          \item Skip if $i{=}m_x{-}1$ or $j{=}m_y{-}1$ or $k{=}m_z{-}1$.
          \item \textbf{$\xi$-face} (guard $j\neq 0 \land k\neq 0$):
                compute metric $g^{11},g^{21},g^{31}$ from \texttt{icsi/ieta/izet};
                call \texttt{Compute\_dscalar\_i} and
                \texttt{getFace2Cell4StencilCsi};
                write \texttt{div[k][j][i].x} and \texttt{visc[k][j][i].x}.
          \item \textbf{$\eta$-face} (guard $i\neq 0 \land k\neq 0$):
                analogous with \texttt{jcsi/jeta/jzet}, writing \texttt{.y}.
          \item \textbf{$\zeta$-face} (guard $i\neq 0 \land j\neq 0$):
                analogous with \texttt{kcsi/keta/kzet}, writing \texttt{.z}.
        \end{itemize}
  \item Halo exchange (\texttt{DMLocalToLocal}) on \texttt{lDivT} and \texttt{lViscT}.
  \item \textbf{Accumulation loop} over the local-with-ghost range:
        \[
          \texttt{rhs}[k][j][i]
          = \texttt{scale}\times
            \bigl[
              (d_x + d_y + d_z)_{\text{conv}}\cdot\mathbf{1}_{f\neq 2}
            + (v_x + v_y + v_z)_{\text{visc}}\cdot\mathbf{1}_{f\neq 1}
            \bigr]
            \times J_{kji},
        \]
        where subscripts $f\neq 2$ / $f\neq 1$ reflect \texttt{formMode}.
\end{enumerate}

\subsection*{Loop fusion rationale}

Prior to this optimisation \texttt{FormT} contained three identical
\texttt{for k for j for i} loops (one per face direction), each traversing
$N_k\times N_j\times N_i$ cells.  Fusing them into a single loop with
per-direction \texttt{if}-guards:
\begin{itemize}
  \item reduces loop-overhead by $3\times$,
  \item improves cache reuse — all three face contributions for cell $(k,j,i)$
        are computed while that cache line is hot,
  \item matches the existing fusion in \texttt{FormU} (momentum equation).
\end{itemize}
The different boundary exclusions ($j{=}0$/$k{=}0$ for $\xi$, etc.) become
lightweight branch conditions inside the single loop body.

% ----------------------------------------------------------------
\section{Time-Stepping Schemes for the Temperature Equation}
\label{sec:teqn-schemes}
% ----------------------------------------------------------------

The active scheme is selected by \texttt{(teqn\_)}$\to$\texttt{ddtScheme}
(set in \texttt{control.dat}).

\subsection{Forward Euler}
\label{ssec:teqn-fe}

\begin{equation}
  T^{n+1} = T^n + \Delta t\,\mathcal{F}(T^n).
\end{equation}

Implemented directly in \texttt{SolveTEqn} via a single \texttt{FormT} call
(\texttt{formMode}=0) followed by \texttt{VecAXPY}.

\subsection{Runge--Kutta 4 (\texttt{rungeKutta4})}
\label{ssec:teqn-rk4}

Classical four-stage explicit RK:
\begin{align}
  k_1 &= \Delta t\,\mathcal{F}(T^n), \notag\\
  k_2 &= \Delta t\,\mathcal{F}(T^n + \tfrac{1}{2}k_1), \notag\\
  k_3 &= \Delta t\,\mathcal{F}(T^n + \tfrac{1}{2}k_2), \notag\\
  k_4 &= \Delta t\,\mathcal{F}(T^n + k_3), \notag\\
  T^{n+1} &= T^n + \tfrac{1}{6}(k_1 + 2k_2 + 2k_3 + k_4). \notag
\end{align}

Stability limit: $\text{CFL} \lesssim 2\sqrt{2}$ (convection-dominated).

\subsection{Backward Euler / SNES (\texttt{backwardEuler})}
\label{ssec:teqn-be}

\subsubsection*{Residual form}

The nonlinear system to solve at each time step is
\begin{equation}
  \mathbf{F}(T^{n+1}) \;=\;
  T^{n+1} - T^n
  - \Delta t\,\mathcal{F}(T^{n+1})
  - \Delta t\,\mathcal{D}(T^{n+1})
  - s_T
  \;=\; \mathbf{0},
  \label{eq:teqn-snes-residual}
\end{equation}
where $\mathcal{D}$ denotes the Rayleigh damping source and $s_T$ the remaining
source terms.
Note: \texttt{sourceT} is added \emph{without} a $\Delta t$ factor (the
convention is that \texttt{sourceT} already encodes a time-integrated
contribution for one step).

\subsubsection*{Explicit Euler predictor}

Before handing the nonlinear system to SNES, an explicit Euler step provides a
warm start:
\begin{equation}
  T^{*} = T^n + \Delta t\,\mathcal{F}(T^n),
  \label{eq:teqn-predictor}
\end{equation}
computed via \texttt{FormT(teqn, RhsTmp, dt, 0)}.
The predictor is stored in \texttt{TmprtTmp} and used as the initial guess for
SNES (\texttt{SNESSolve(snes, NULL, TmprtTmp)}).

\subsubsection*{SNES configuration}

\begin{center}
\begin{tabular}{ll}
\toprule
Parameter & Value / source \\
\midrule
Type & \texttt{SNESNEWTONLS} \\
Relative tolerance & \texttt{teqn->solverTol} \\
Max iterations & \texttt{teqn->solverItMax} \\
KSP inner solver & inherited from PETSc defaults \\
\bottomrule
\end{tabular}
\end{center}

\subsection{BDF2 / SNES (\texttt{BDF2})}
\label{ssec:teqn-bdf2}

\subsubsection*{Motivation}

Backward Euler is only first-order accurate in time; Crank--Nicolson is second-order but neutrally stable ($|G|=1$ at all wavenumbers), meaning any
numerical perturbation is neither damped nor amplified.  In practice the neutral
stability of Crank--Nicolson leads to divergence in stratified LES because the
LES sub-grid model, the buoyancy coupling, and SNES residual tolerances all
deposit small high-frequency errors that accumulate without bound.  
BDF2 (second-order Backward Differentiation Formula) provides second-order
accuracy with mild A-stable dissipation ($|G|<1$ at high wavenumbers), which
damps these errors without corrupting the resolved low-wavenumber dynamics.

\subsubsection*{Residual form}

BDF2 uses a three-level time-derivative approximation.  At step $n+1$ the
residual to solve is
\begin{equation}
  \mathbf{F}(T^{n+1}) \;=\;
  \frac{3}{2}\,T^{n+1} - 2\,T^n + \frac{1}{2}\,T^{n-1}
  - \Delta t\,\mathcal{F}(T^{n+1})
  - s_T
  \;=\; \mathbf{0},
  \label{eq:teqn-bdf2-residual}
\end{equation}
where $\mathcal{F}(T^{n+1})$ is the full right-hand side (convection + diffusion
+ Rayleigh damping) evaluated at the new time level.  On the first step
($\text{it} = \text{itStart}$) BDF2 falls back to Backward Euler (only $T^n$
available), giving first-order accuracy for that step only.

Implementation in \texttt{TeqnSNES}: when \texttt{ddtScheme == "BDF2"} and
$\text{it} > \text{itStart}$ the time-derivative terms in the residual vector
are assembled as
\begin{lstlisting}[language=C,basicstyle=\ttfamily\small]
VecAXPY(Rhs, -3.0/2.0, T);           // -3/2 T^{n+1}  (from SNES iterate)
VecAXPY(Rhs,  2.0,     teqn->Tmprt_o);   //  +2   T^n
VecAXPY(Rhs, -1.0/2.0, teqn->Tmprt_oo); //  -1/2 T^{n-1}
\end{lstlisting}
\noindent
For the first step the existing two-level fall-back ($-T^{n+1}+T^n$, i.e.\ Backward Euler)
is applied automatically.

\subsubsection*{History rotation}

At the beginning of each time step (before the step is taken):
\begin{lstlisting}[language=C,basicstyle=\ttfamily\small]
// main.c — start of time step
if(ddtScheme == "BDF2")
    VecCopy(Tmprt_o, Tmprt_oo);   // T^{n-1} <- T^n
VecCopy(Tmprt, Tmprt_o);          // T^n     <- T^{current}
\end{lstlisting}
\noindent
The additional storage vector \texttt{Tmprt\_oo} is allocated in
\texttt{InitializeTEqn} only when \texttt{ddtScheme == "BDF2"}.

\subsubsection*{Ghost-cell behaviour}

At ghost/IBM cells \texttt{resetNonResolvedCellCentersScalar} zeros the
residual before the time-derivative terms are added.  The BDF2
residual at a ghost cell is therefore
$-\tfrac{3}{2}T^{n+1}_{\text{ghost}} + 2T^n_{\text{ghost}} - \tfrac{1}{2}T^{n-1}_{\text{ghost}} = 0$,
which effectively extrapolates the ghost value as
$T^{n+1}_{\text{ghost}} = \tfrac{4}{3}T^n - \tfrac{1}{3}T^{n-1}$.
For time-invariant boundary conditions ($T^n_{\text{ghost}}=T^{n-1}_{\text{ghost}}=T_{\text{BC}}$)
this extrapolation recovers $T_{\text{BC}}$ exactly; for time-varying BCs the
error is $O(\Delta t^2)$, consistent with the global second-order accuracy.

\subsubsection*{SNES configuration}

Identical to Backward Euler: \texttt{SNESNEWTONLS}, matrix-free MFFD Jacobian
(\texttt{JT}), BiCGStab or GMRES inner KSP (\texttt{ksp}/\texttt{pc}),
tolerance controlled by \texttt{-absTolT} and \texttt{-relTolT}.

\subsection{IMEX-CNAB-BE (\texttt{IMEX})}
\label{ssec:teqn-imex}

\subsubsection*{Motivation}

Treating the nonlinear convection explicitly avoids a nonlinear solve;
treating diffusion implicitly removes the thermal diffusivity constraint on
$\Delta t$.  Unlike the momentum IMEX which uses Crank--Nicolson (CN) for
diffusion, the temperature IMEX uses \emph{backward Euler} for diffusion.
This choice decouples $T^{n+1}$ from $\mathbf{u}^{n+1}$: the velocity field
is not yet updated when the temperature KSP is called, so CN would require
an estimate of $\nabla\!\cdot\!(\kappa_{\text{eff}}\nabla T)$ at $t^{n+1}$
using a velocity that has not yet converged.

\subsubsection*{Operator splitting}

Write $\mathcal{F} = \mathcal{N} + \mathcal{D}$ (convection + diffusion).
The IMEX-CNAB-BE update is
\begin{equation}
  \underbrace{(I - \Delta t\,\mathcal{D})}_{A}\,T^{n+1}
  = T^n + \Delta t\,b_T,
  \label{eq:teqn-imex-system}
\end{equation}
with the explicit right-hand side
\begin{equation}
  b_T = c_1\,\mathcal{N}(T^n) + c_2\,\mathcal{N}(T^{n-1})
       + \mathcal{D}_{\text{damp}}(T^n) + s_T / \Delta t,
  \label{eq:bT}
\end{equation}
where $c_1, c_2$ are the Adams--Bashforth~2 blending coefficients
(first step: $c_1{=}1, c_2{=}0$; subsequent: $c_1{=}3/2, c_2{=}{-}1/2$)
and $\mathcal{D}_{\text{damp}}$ is the Rayleigh damping operator.

\subsubsection*{Matrix-vector product — \texttt{IMEXTMatVec}}

The linear operator $A = I - \Delta t\,\mathcal{D}$ is represented as a
MatShell.  Its action on a vector $v$ is:

\begin{equation}
  A\,v = v - \Delta t\,\mathcal{D}(v)
       = v - \Delta t\,\texttt{FormT}(v,\,\texttt{scale}{=}1.0,\,\texttt{DIFF\_ONLY}).
  \label{eq:IMEXTMatVec}
\end{equation}

The \texttt{scale} argument is always $1.0$; the factor $\Delta t$ is applied
externally in \texttt{IMEXTMatVec} before subtraction.
This is in contrast to the momentum \texttt{IMEXMatVec}, which uses a
Crank--Nicolson factor of $0.5$.

\begin{tcolorbox}[colback=blue!5,colframe=blue!50!black,title=Key difference from momentum IMEX]
  \textbf{Momentum (IMEX-CNAB):}
  $A_U\,v = v - \tfrac{\Delta t}{2}\,\mathcal{D}_U(v)$
  \hfill (Crank--Nicolson, $\theta=0.5$)\\[4pt]
  \textbf{Temperature (IMEX-BE):}
  $A_T\,v = v - \Delta t\,\mathcal{D}_T(v)$
  \hfill (Backward Euler, $\theta=1.0$)\\[4pt]
  The backward-Euler choice avoids requiring $\mathbf{u}^{n+1}$ for
  the diffusion term when $\kappa_{\text{eff}}$ depends on SGS quantities
  computed from velocity.
\end{tcolorbox}

\subsubsection*{\texttt{TeqnIMEX} — seven-step procedure}

\begin{enumerate}
  \item[\textbf{1.}] Compute $\mathcal{N}(T^n)$:
        \texttt{FormT(teqn, conv, 1.0, CONV\_ONLY)}.
  \item[\textbf{2.}] Build $b_T = c_1\,\mathcal{N}^n + c_2\,\mathcal{N}^{n-1}
        + \mathcal{D}_{\text{damp}} + s_T/\Delta t$
        and store in \texttt{teqn->bT}.
  \item[\textbf{3.}] Form the RHS of the linear system:
        \texttt{TmprtTmp = $T^n$ + $\Delta t\,b_T$} \;
        (i.e.\ \texttt{VecWAXPY(TmprtTmp, dt, bT, Tmprt)})
        then add $s_T$ (\texttt{sourceT(teqn, TmprtTmp, 1.0)}).
  \item[\textbf{4.}] Set up the MatShell $A$ (\texttt{MatCreateShell},
        \texttt{MatShellSetOperation}) pointing to \texttt{IMEXTMatVec}.
  \item[\textbf{5.}] Configure KSP: type from \texttt{kspType} (default
        \texttt{KSPBCGSL}), tolerances from \texttt{solverTol} /
        \texttt{solverItMax}.
  \item[\textbf{6.}] Solve $A\,T^{n+1} = \texttt{TmprtTmp}$.
  \item[\textbf{7.}] Store $\mathcal{N}^n \to \mathcal{N}^{n-1}$
        (\texttt{VecCopy(RhsConv, RhsConv\_o)}).
\end{enumerate}

\subsubsection*{Control parameters for IMEX}

\begin{center}
\begin{tabular}{lll}
\toprule
\textbf{Parameter} & \textbf{\texttt{teqn\_} field} & \textbf{Default / notes} \\
\midrule
KSP solver type & \texttt{kspType} &
  \texttt{"bcgsl"} (BiCGStab$(\ell{=}2)$); also accepts
  \texttt{"gmres"}, \texttt{"fgmres"} \\
GMRES restart & \texttt{gmresRestart} & 30 \\
Relative tolerance & \texttt{solverTol} & $10^{-5}$ \\
Max iterations & \texttt{solverItMax} & 200 \\
\bottomrule
\end{tabular}
\end{center}

\subsubsection*{Buffer fields}

\begin{center}
\begin{tabular}{lll}
\toprule
\textbf{Field} & \textbf{Type} & \textbf{Content} \\
\midrule
\texttt{bT}       & \texttt{Vec} & Explicit RHS blended convection + damping + source \\
\texttt{RhsConv}  & \texttt{Vec} & $\mathcal{N}(T^n)$ — current convective RHS \\
\texttt{RhsConv\_o} & \texttt{Vec} & $\mathcal{N}(T^{n-1})$ — previous convective RHS \\
\texttt{JvIMEX}   & \texttt{Mat} & MatShell for $A = I - \Delta t\,\mathcal{D}$ \\
\texttt{kspIMEX}  & \texttt{KSP} & PETSc KSP context for IMEX linear solve \\
\bottomrule
\end{tabular}
\end{center}

% ----------------------------------------------------------------
\section{U--T Buoyancy Coupling and the T Predictor}
\label{sec:teqn-coupling}
% ----------------------------------------------------------------

\subsection*{The coupling problem}

The main time loop advances fields in the following order:

\begin{center}
\texttt{Buoyancy}$(T^n)$ $\;\to\;$
\texttt{SolveUEqn} $\;\to\;$
\texttt{SolvePEqn} $\;\to\;$
\texttt{SolveTEqn}
\end{center}

\texttt{Buoyancy} writes the buoyancy acceleration $\mathbf{b}(T^n)$ into
\texttt{bTheta}, which \texttt{FormU} then uses.  Because $T^n$ is one step
behind, the buoyancy seen by the momentum update is lagged:

\begin{equation}
  \mathbf{u}^{n+1}
  = \text{advance}\bigl(\mathbf{u}^n,\,\mathbf{b}(T^n)\bigr),
  \quad\text{error }\sim \Delta t\,\partial_t\mathbf{b} = O(\Delta t).
\end{equation}

For stably-stratified flow this $O(\Delta t)$ phase error between velocity
and temperature introduces a spurious phase lag in the buoyancy restoring
force.  When the Brunt--Väisälä period $\tau_N = 2\pi/N$ is comparable to
$\Delta t$, the lag grows into a numerical instability.

\subsection*{T predictor (\texttt{TmprtPredictor})}

To reduce the coupling error to $O(\Delta t^2)$, the momentum step can
optionally be preceded by a single cheap convection-only forward-Euler step:

\begin{equation}
  T^* = T^n + \Delta t\,\mathcal{N}(T^n),
  \label{eq:Tpredictor}
\end{equation}

where $\mathcal{N}$ is the convective operator evaluated with the
\emph{already-updated} $\mathbf{u}^{n+1}$ contravariant fluxes
(formMode~=~1, scale~=~1.0).
The modified time loop becomes:

\begin{center}
\texttt{TmprtPredictor}: $T^*=T^n+\Delta t\,\mathcal{N}(T^n)$
$\;\to\;$
\texttt{Buoyancy}$(T^*)$
$\;\to\;$
\texttt{SolveUEqn}
$\;\to\;$
\texttt{SolvePEqn}
$\;\to\;$
\texttt{TmprtRestoreFromOld}: $\texttt{lTmprt}\leftarrow T^n$
$\;\to\;$
\texttt{SolveTEqn}
\end{center}

Since $T^* \approx T^{n+1/2}$ to first order,
the buoyancy error is reduced to $O(\Delta t^2)$:

\begin{equation}
  \mathbf{b}(T^*) = \mathbf{b}(T^n) + \Delta t\,\partial_t\mathbf{b}\big|_n
  + O(\Delta t^2),
\end{equation}

which matches the second-order accuracy of the velocity schemes (RK4, IMEX,
backward Euler).

\textbf{Why convection only?}
Advection is the dominant mechanism that repositions the temperature inversion
layer within one time step.
Diffusion changes $T$ on the diffusion time-scale
$\tau_\kappa = h^2/\kappa_{\text{eff}} \gg \Delta t$ for large CFL, so
neglecting it in the predictor introduces an error well below $O(\Delta t^2)$.
The predictor therefore adds only one \texttt{FormT(CONV\_ONLY)} call
(roughly one-fifth of a full \texttt{FormT}) and one halo scatter per step.

\textbf{Restore before \texttt{SolveTEqn}.}
The predictor temporarily stores $T^*$ in \texttt{Tmprt}/\texttt{lTmprt}.
Before \texttt{SolveTEqn} is called, \texttt{TmprtRestoreFromOld} copies
\texttt{Tmprt\_o} (which holds $T^n$, saved at the start of the time step)
back into \texttt{Tmprt}/\texttt{lTmprt}.  This ensures that \texttt{TeqnIMEX}
Step~2 (which reads \texttt{lTmprt} as $T^n$) and \texttt{TeqnSNES} (which
uses \texttt{Tmprt\_o} as the old state) are both consistent: the T solver
always advances from $T^n$, not from $T^*$.

\subsection*{Activation}

The predictor is controlled by the \texttt{-teqnPredictorU} PETSc option
(set in \texttt{control.dat} or on the command line):

\begin{lstlisting}[language={}]
-teqnPredictorU    forwardEuler;    # or "none" (default)
\end{lstlisting}

The option is read in \texttt{InitializeUEqn} and stored in
\texttt{ueqn\_::teqnPredictorScheme}.  It is independent of both the U and T
\texttt{ddtScheme} choices (it applies to any combination).

% ----------------------------------------------------------------
\section{Key Design Decisions}
\label{sec:teqn-design}
% ----------------------------------------------------------------

\begin{enumerate}
  \item \textbf{Backward Euler for diffusion (not CN)}: Avoids coupling
        $T^{n+1}$ to $\kappa_{\text{eff}}(\mathbf{u}^{n+1})$ before the
        velocity update; the linear system for $T$ is self-contained.

  \item \textbf{\texttt{sourceT} convention}: The source term is added
        \emph{without} scaling by $\Delta t$ in both \texttt{TeqnSNES} and
        \texttt{TeqnIMEX}.  In \texttt{TeqnSNES} the \texttt{Rhs} vector is
        first scaled by $\Delta t$ with \texttt{VecScale(Rhs, dt)}, then
        \texttt{sourceT(teqn, Rhs, 1.0)} is called; the source therefore
        enters at its physical magnitude, not scaled by $\Delta t$.
        In \texttt{TeqnIMEX} the source is added to the already-built
        right-hand side \texttt{TmprtTmp = $T^n + \Delta t\,b_T$}, again
        without an extra $\Delta t$ factor.

  \item \textbf{Explicit predictor for SNES warm start}: Used as an initial
        guess for the SNES (\texttt{backwardEuler}) to reduce nonlinear
        iterations; distinct from the buoyancy predictor above.

  \item \textbf{Loop fusion in \texttt{FormT}}: The original three separate
        triple-nested loops (one per face direction) were replaced by a single
        fused loop with per-direction guards, matching the optimisation
        already present in \texttt{FormU}.

  \item \textbf{T predictor before \texttt{SolveUEqn}}: Optional
        (\texttt{-teqnPredictorU forwardEuler}) — reduces U--T buoyancy
        coupling error from $O(\Delta t)$ to $O(\Delta t^2)$ at the cost
        of one extra \texttt{FormT(CONV\_ONLY)} call per time step.
\end{enumerate}

% ================================================================
\section{Pressure Boundary-Condition Ordering Bug and Fix}
\label{sec:peqn-bc-bug}
% ================================================================

\subsection{Background: pressure update pipeline}
\label{sec:peqn-bc-bug-bg}

At the end of each pressure-projection step the following sequence is
executed inside \texttt{SolvePEqn} (after the Poisson solver converges):

\begin{enumerate}
  \item \textbf{\texttt{phiToPhi}} — converts the raw solver output
        ($\phi_2$) to the actual pressure correction $\phi$.
  \item \textbf{\texttt{UpdatePressure}} — adds $\phi$ to $p$ for every
        fluid cell, computes and subtracts the domain mean (compatibility
        condition), then scatters $P \to \mathtt{lP}$.
  \item \textbf{\texttt{SetPressureReference}} — pins the gauge by
        subtracting $p_{1,1,1}$ from every non-solid cell, then
        scatters $P \to \mathtt{lP}$.
  \item \textbf{\texttt{UpdatePressureBCs}} — reads \texttt{lP} and
        reflects the correct Neumann (zero-gradient) or specified
        boundary values into the ghost rows of $P$, then scatters
        $P \to \mathtt{lP}$.
  \item \textbf{\texttt{ProjectVelocity}} — uses \texttt{lP} to correct
        the contravariant fluxes $U^*_{\text{cont}} \to U_{\text{cont}}$.
  \item \textbf{Next step}: \texttt{GradP} builds the pressure-gradient
        forcing from \texttt{lP} for the momentum equation.
\end{enumerate}

\subsection{The bug: premature call to \texttt{UpdatePressureBCs}}
\label{sec:peqn-bc-bug-desc}

In the original code \texttt{UpdatePressureBCs} was called from
\emph{inside} \texttt{UpdatePressure}, between the first accumulation
loop (which adds $\phi$ to $p$ for fluid cells) and the
\texttt{MPI\_Allreduce} that computes the domain mean.  The full
execution order was therefore:

\begin{enumerate}
  \item Fluid cells: $p \mathrel{+}= \phi$
  \item \textbf{\texttt{UpdatePressureBCs}} — sets ghost-row values of
        $P$ from \texttt{lP}, which still holds data from the
        \emph{previous} time step at this point.
  \item \texttt{MPI\_Allreduce} — compute mean $\bar{p}$
  \item Fluid cells: $p \mathrel{-}= \bar{p}$
  \item \textbf{All remaining cells} (\texttt{else} branch):
        $p_{k,j,i} = 0$ — this branch covers ghost rows as well as
        IBM/overset cells, so it \emph{overwrites} the values that
        \texttt{UpdatePressureBCs} had just written.
  \item Scatter $P \to \mathtt{lP}$: ghost cells in \texttt{lP} become
        exactly $0$.
  \item \texttt{SetPressureReference}: subtracts the gauge shift
        $g = p_{1,1,1}$ from \emph{all} non-IBM cells, including ghost
        rows $\Rightarrow$ ghost cells in \texttt{lP} become $-g$.
\end{enumerate}

After this sequence every ghost cell in \texttt{lP} holds $-g$ instead
of the correct zero-gradient extrapolation.

\subsection{Consequences}
\label{sec:peqn-bc-bug-consequences}

\paragraph{Spurious pressure gradient at boundaries.}
\texttt{GradP} evaluates face-centred pressure gradients by differencing
values of \texttt{lP} at neighbouring cell centres.  For a cell adjacent
to any non-periodic boundary the stencil reads one interior value $p_I$
and one ghost value $p_g$.  With the bug, $p_g = -g$, so

\begin{equation}
  \nabla p\big|_{\partial\Omega} \;\approx\; \frac{p_I - (-g)}{h}
  = \frac{p_I + g}{h},
\end{equation}

an error of magnitude $g/h$ that is \emph{wrong by construction}
and changes sign and amplitude with the time-varying gauge shift $g$.

\paragraph{Contaminated Poisson right-hand side.}
The spurious boundary gradient enters the momentum equation and
corrupts $\mathbf{U}^{n+1}_{\text{cont}}$ at all boundary-adjacent
faces.  \texttt{SetRHS} computes the divergence of
$\mathbf{U}^{n+1}_{\text{cont}}$, so the Poisson right-hand side
$\nabla \cdot \mathbf{U}^*/\Delta t$ is also contaminated, even though
the matrix stencil itself is correct.  The solver faithfully minimises
the residual of the \emph{wrong} system.

\paragraph{Closed feedback loop and pressure oscillations.}
Because $g$ is the domain-mean pressure, which changes every step due
to open-boundary flux imbalances, buoyancy, or mean-gradient forcing,
the magnitude and sign of the error varies in time.  This produces a
self-sustaining oscillation whose amplitude can grow to the point where
the instantaneous min/max pressure range is far larger than the
physically meaningful variation; subtracting the spatial mean does not
suppress this oscillation because the contamination is spatially
structured (concentrated at boundaries and propagated inward by the
elliptic solver), not a uniform DC shift.

\subsection{The fix}
\label{sec:peqn-bc-bug-fix}

\texttt{UpdatePressureBCs} is removed from \texttt{UpdatePressure} and
called in \texttt{SolvePEqn} \emph{after} \texttt{SetPressureReference},
so that it reads the fully corrected, gauge-fixed \texttt{lP}:

\begin{lstlisting}[caption={Corrected call order in \texttt{SolvePEqn}
  (\texttt{src/peqn.c}).},label=lst:peqn-fix]
phiToPhi(peqn);
UpdatePressure(peqn);          // p += phi, mean-subtracted, P->lP
SetPressureReference(peqn);    // gauge shift, P->lP
UpdatePressureBCs(peqn);       // ghost BCs from correct lP, P->lP  <-- moved here
ProjectVelocity(peqn);         // correct fluxes using final lP
\end{lstlisting}

Inside \texttt{UpdatePressure} the call to \texttt{UpdatePressureBCs}
is deleted.  \texttt{lP} is now consistent at every point where it is
subsequently read (\texttt{ProjectVelocity}, \texttt{GradP}).

\begin{warnbox}[Affected solvers]
The bug was present in all time-stepping schemes because
\texttt{SolvePEqn} is called identically from the Forward~Euler,
Runge--Kutta~4, Crank--Nicolson, and IMEX loops.  All schemes benefit
from the reordering.
\end{warnbox}

% ============================================================
\section{Pressure Poisson Solver (\texttt{peqn})}
\label{sec:peqn}

TOSCA uses a fractional-step (operator-splitting) projection method.
After the momentum equation advances the velocity to a predicted state
$\widetilde{U}^*$ that is not necessarily divergence-free, a pressure
correction $\phi$ is obtained by solving a Poisson equation; the
corrected contravariant flux is then projected onto the divergence-free
manifold.  The entire sequence is orchestrated by \texttt{SolvePEqn}.

% ------------------------------------------------------------
\subsection{Overall Sequence (\texttt{SolvePEqn})}
\label{sec:peqn-sequence}

\begin{lstlisting}[caption={Body of \texttt{SolvePEqn} (\texttt{src/peqn.c}).},
                   label=lst:solvePEqn]
// 1. IBM flux correction (IBM/overset only)
AdjustIBMFlux(peqn);

// 2. Dynamic-IBM matrix refresh (if needed)
//    HYPRE: in-place re-initialise (avoids dangling-pointer hang)
//    PETSc: destroy/recreate KSP
SetCoeffMatrix(peqn);

// 3. Divergence-based RHS + compatibility condition
SetRHS(peqn);

// 4. Linear solve  (HYPRE or PETSc, see Sec. sec:peqn-solvers)
//    Initial guess = phi from previous timestep (warm start)
// -> overwrites peqn->phi

// 5. Post-solve: subtract mean of phi (null-space removal)
SubtractAverageHypre / SubtractAveragePETSc(peqn);

// 6. Scatter flat phi -> 3-D DMDA field
phiToPhi(peqn);

// 7. Accumulate pressure and subtract its mean
UpdatePressure(peqn);         // p^{n+1} = p^n + Phi, mean-subtracted

// 8. Gauge fix: p[1][1][1] = 0
SetPressureReference(peqn);

// 9. Fill ghost cells from final lP
UpdatePressureBCs(peqn);      // MUST be after SetPressureReference

// 10. Project contravariant fluxes
ProjectVelocity(peqn);
\end{lstlisting}

Steps 7--9 must appear in the order shown; see
Section~\ref{sec:peqn-bc-bug} for the consequences of calling
\texttt{UpdatePressureBCs} earlier.

% ------------------------------------------------------------
\subsection{DOF Numbering (\texttt{SetPoissonConnectivity})}
\label{sec:peqn-connectivity}

\texttt{SetPoissonConnectivity} assigns a global degree-of-freedom
index to every interior cell (including IBM and overset cells).  The
mapping is:

\begin{equation}
  \texttt{matID}(i,j,k) = \texttt{lid}[k][j][i] + \texttt{nUntilThisRank},
\end{equation}

where \texttt{lid} is a rank-local counter that visits cells in
$(k,j,i)$ order, and \texttt{nUntilThisRank} is the accumulated DOF
count of all lower-ranked processes.  The per-rank range
$[\texttt{thisRankStart},\texttt{thisRankEnd}]$ is communicated via
\texttt{MPI\_Allreduce} and used by both the HYPRE and PETSc paths for
consistent parallel matrix layout.

IBM/overset cells receive a valid DOF ID but their rows in the matrix
are treated as trivial equations $\phi_i = 0$ (see below).  The
\texttt{phi} PETSc vector is sized to \texttt{gndof[rank]} local
entries, matching the DOF layout exactly.

% ------------------------------------------------------------
\subsection{Poisson Matrix (\texttt{SetCoeffMatrix})}
\label{sec:peqn-matrix}

The discrete Laplacian on the curvilinear grid gives a 19-point stencil
per fluid cell.  At each cell-face $\alpha\in\{i,j,k\}$ the contribution
involves the metric-tensor component

\begin{equation}
  g^{(\alpha)}_{\beta\gamma}
  = \bigl(\nabla\xi^\beta \cdot \nabla\xi^\gamma\bigr)\,J_\alpha,
\end{equation}

where $J_\alpha$ is the face Jacobian (\texttt{iaj}/\texttt{jaj}/\texttt{kaj}).
Each face contributes up to nine off-diagonal entries (the
$3\times3$ cross-derivative block at the face), leading to the 19-point
stencil (1 diagonal + 6 face-centred + 12 edge-centred).

IBM and overset cells receive the trivial row $A_{ii}=1$,
$A_{ij}=0$~($j\neq i$), so their solution is identically zero and they
contribute nothing to the divergence-constraint.

A face contribution is \emph{excluded} whenever the neighbour cell in
the direction of differentiation is an IBM cell or a domain boundary
under a zero-gradient (Neumann) condition; this is equivalent to setting
the normal derivative of $\phi$ to zero at that face, which is the
correct Neumann boundary condition.

% ------------------------------------------------------------
\subsection{Right-Hand Side and Compatibility (\texttt{SetRHS})}
\label{sec:peqn-rhs}

The RHS of the Poisson system is the divergence of the predicted
contravariant flux divided by $\Delta t$:

\begin{equation}
  b_{ijk}
  = \frac{1}{\Delta t}
    \bigl(
      \widetilde{U}^{\xi}_{i+\tfrac12} - \widetilde{U}^{\xi}_{i-\tfrac12}
    + \widetilde{U}^{\eta}_{j+\tfrac12} - \widetilde{U}^{\eta}_{j-\tfrac12}
    + \widetilde{U}^{\zeta}_{k+\tfrac12} - \widetilde{U}^{\zeta}_{k-\tfrac12}
    \bigr).
\end{equation}

IBM/overset cells receive $b_{ijk}=0$ (consistent with $A_{ii}=1$,
solution $\phi=0$).

The Poisson system has a constant null space (pure Neumann problem),
so the compatibility condition $\sum_\Omega b_i = 0$ must be enforced.
\texttt{SetRHS} does this by computing the mean divergence over fluid
(non-IBM, non-overset) cells via \texttt{MPI\_Allreduce} and
subtracting it from every fluid-cell RHS entry.

% ------------------------------------------------------------
\subsection{Null-Space Removal from the Solution}
\label{sec:peqn-nullspace}

After each linear solve the constant mode in $\phi$ is removed by
\texttt{SubtractAverageHypre} / \texttt{SubtractAveragePETSc}: the
mean of $\phi$ is computed over \emph{fluid and calculated cells only}
(filtered by \texttt{isFluidCell} \&\& \texttt{isCalculatedCell}) and
subtracted from those same cells.  The two functions are numerically
equivalent after the fix described in Section~\ref{sec:peqn-bc-bug}.

IBM/overset cells, whose solution is identically zero, are excluded from
both the mean computation and the subtraction; including them would
dilute the mean and introduce a spurious shift.

% ------------------------------------------------------------
\subsection{Linear Solvers}
\label{sec:peqn-solvers}

The solver backend is selected via \texttt{solverType} in
\texttt{control.dat}.

\subsubsection{HYPRE path (\texttt{solverType = HYPRE})}

The outer Krylov method is either GMRES (type~1) or PCG (type~2), both
wrapped around a BoomerAMG preconditioner configured in
\texttt{CreateHypreSolver}:

\begin{lstlisting}[caption={BoomerAMG configuration in \texttt{CreateHypreSolver}.}]
HYPRE_BoomerAMGSetInterpType(hyprePC, 13);    // FF1 interpolation
HYPRE_BoomerAMGSetStrongThreshold(hyprePC, amgThresh); // 0.5--0.6
HYPRE_BoomerAMGSetCoarsenType(hyprePC, amgCoarsenType);// e.g. HMIS=10
HYPRE_BoomerAMGSetCycleType(hyprePC, 1);      // V-cycle
HYPRE_BoomerAMGSetMaxIter(hyprePC, 3);        // 3 AMG sweeps per KSP iter
HYPRE_BoomerAMGSetSmoothType(hyprePC, 7);     // Pilut smoother
HYPRE_BoomerAMGSetSmoothNumSweeps(hyprePC, 3);
\end{lstlisting}

BoomerAMG is $\mathcal{O}(N)$ scalable; convergence is independent of
grid size when the strong threshold and coarsening type are correctly
tuned (0.5 for Cartesian-like meshes, 0.6 for highly distorted ones).

The BoomerAMG hierarchy is built (setup) only at the first timestep or
when the IBM body moves.  On subsequent steps the existing hierarchy is
reused, making each solve essentially free of setup cost.

\paragraph{Dynamic IBM.}
When the IBM body moves, the matrix values change but the DOF layout
does not (\texttt{SetPoissonConnectivity} always assigns IDs to all
interior cells regardless of IBM status).  The matrix and vectors are
therefore \emph{re-initialised in place} via
\texttt{HYPRE\_IJMatrixInitialize} and \texttt{HYPRE\_IJVectorInitialize}
rather than destroyed and recreated:

\begin{warnbox}[Dynamic-IBM HYPRE pattern]
\texttt{HYPRE\_IJMatrixDestroy} frees the underlying ParCSR object.
However, \texttt{HYPRE\_BoomerAMGSetup} retains an internal pointer to
that object; destroying it leaves a dangling pointer that causes a
random use-after-free hang on the next \texttt{HYPRE\_BoomerAMGSetup}
call.  Re-initialising in place keeps the ParCSR pointer valid and
eliminates the hang.
\end{warnbox}

\subsubsection{PETSc path (\texttt{solverType = PETSc})}

The PETSc path uses FGMRES with a PCGAMG algebraic-multigrid
preconditioner, configured in \texttt{CreatePETScSolver}:

\begin{lstlisting}[caption={PETSc solver configuration in \texttt{CreatePETScSolver}.}]
PCSetType(petscPC, PCGAMG);
PCGAMGSetNSmooths(petscPC, 1);      // 1 smoother sweep per level
KSPSetPCSide(ksp, PC_RIGHT);        // FGMRES only supports right preconditioning
KSPSetType(ksp, KSPFGMRES);         // FGMRES: supports variable PC
KSPGMRESSetRestart(ksp, 20);        // restart after 20 vectors
KSPSetInitialGuessNonzero(ksp, PETSC_TRUE); // warm start
KSPSetUp(ksp);                      // build AMG hierarchy
\end{lstlisting}

Key design choices:
\begin{itemize}
  \item \textbf{FGMRES} (not PGMRES): AMG preconditioners can vary
    between applications (e.g.\ after the hierarchy is re-built);
    FGMRES handles variable preconditioners correctly, PGMRES assumes
    the preconditioner is fixed.
  \item \textbf{Restart 20}: PCGAMG typically converges in fewer than
    20 iterations for Poisson on well-resolved LES grids; a restart of
    100 (the previous default) pre-allocated 100 basis vectors
    needlessly.
  \item \textbf{Warm start}: \texttt{KSPSetInitialGuessNonzero} tells
    KSP to use the current content of \texttt{phi} as initial guess,
    carrying over the solution from the previous timestep.
\end{itemize}

For dynamic IBM the PETSc path performs a full
\texttt{DestroyPETScSolver}+\texttt{CreatePETScSolver} cycle, which
triggers a PCGAMG hierarchy rebuild via \texttt{KSPSetUp}.  This is
more expensive than the HYPRE in-place reuse but is safe
(\texttt{KSPDestroy} is a clean teardown with no dangling-pointer risk).

% ------------------------------------------------------------
\subsection{Post-Solve Pressure Update}
\label{sec:peqn-pupdate}

\subsubsection{\texttt{phiToPhi}: flat vector to 3-D field}

The flat PETSc vector \texttt{phi} (DOF-ordered, length
\texttt{thisRankSize}) is scattered into the 3-D global DMDA vector
\texttt{Phi} and its ghost-padded local counterpart \texttt{lPhi}.
IBM and overset cells receive $\Phi = 0$ explicitly.

\subsubsection{\texttt{UpdatePressure}: accumulate and de-mean}

\begin{equation}
  p^{n+1}_{ijk} = p^n_{ijk} + \Phi_{ijk}
  \quad \text{for fluid cells},
\end{equation}

followed by subtraction of the spatial mean of $p^{n+1}$ over all
fluid-and-calculated cells (via \texttt{MPI\_Allreduce}).  This
intermediate de-meaning prevents unbounded pressure drift but does not
constitute the final gauge fix.

\subsubsection{\texttt{SetPressureReference}: gauge fix}

The cell $(i,j,k)=(1,1,1)$ is chosen as pressure reference.  Its
value is broadcast via \texttt{MPI\_Allreduce} (the rank that owns
the cell loads it into the local contribution; all other ranks
contribute zero) and then subtracted from every non-IBM,
non-zeroed cell.  After this step $p[1][1][1]=0$ exactly.

\subsubsection{\texttt{ProjectVelocity}: divergence-free projection}

The corrected contravariant flux is

\begin{equation}
  U^{n+1}_{\text{ctv},\alpha} =
  \widetilde{U}^*_{\text{ctv},\alpha}
  - \Delta t\,
    \bigl((\nabla\xi^\beta\cdot\nabla\xi^\alpha)J_\alpha\bigr)
    \frac{\partial\phi}{\partial\xi^\beta},
\end{equation}

where the pressure-gradient difference $\partial\phi/\partial\xi^\beta$
at each face is evaluated using the same stencil as
\texttt{SetCoeffMatrix}, ensuring exact consistency between the
projection and the Poisson operator: the projected fluxes are
divergence-free to machine precision.

% ------------------------------------------------------------
\subsection{Pressure Gradient for the Momentum Equation (\texttt{GradP})}
\label{sec:peqn-gradp}

\texttt{GradP} computes the contravariant pressure-gradient vector
$\mathbf{dP}$ used as a source term in the momentum equation.  It
employs \emph{adaptive order}:

\begin{itemize}
  \item 4th-order central differences in the interior, away from
    IBM/overset interfaces and domain boundaries.
  \item 2nd-order central differences within one cell of an IBM or
    overset interface, or within two cells of a domain boundary.
  \item 1st-order one-sided (upwind) differences at domain-boundary
    corner faces where two non-periodic boundaries meet (to avoid
    referencing out-of-domain ghost cells).
\end{itemize}

The stencil switches are applied per face and per direction
independently, so a face near one boundary can use 4th order
in the two tangential directions while using 2nd order in the
normal direction.

\begin{notebox}[Note: \texttt{GradP} vs.\ \texttt{ProjectVelocity} stencils]
  \texttt{GradP} (momentum source) and \texttt{ProjectVelocity}
  (flux correction) use deliberately \emph{different} stencils for
  $\nabla p$.  \texttt{ProjectVelocity} uses the exact transpose of
  the \texttt{SetCoeffMatrix} stencil (ensuring machine-precision
  divergence-free projection), while \texttt{GradP} uses the
  adaptive-order scheme above for accuracy in the interior.  This
  is intentional and does not introduce inconsistency in the
  time-integration.
\end{notebox}


% ============================================================
%  PART III — OVERSET METHOD
% ============================================================
\clearpage
\part*{Part III — Overset Method}
\addcontentsline{toc}{part}{Part III — Overset Method}

% ----------------------------------------------------------------
\section{Concept and Domain Hierarchy}
\label{sec:overset-concept}
% ----------------------------------------------------------------

TOSCA's overset (chimera) method couples two or more structured meshes that
partially overlap in physical space.
The standard configuration is a \textbf{two-level cascade}:

\begin{center}
\begin{tabular}{lll}
\toprule
\textbf{Role} & \textbf{Alias} & \textbf{Purpose} \\
\midrule
Background (parent) & coarse mesh & large-scale domain; provides far-field BCs \\
Fine insert (child)  & fine mesh   & resolved region (turbine, terrain detail) \\
\bottomrule
\end{tabular}
\end{center}

The fine mesh is fully embedded inside the background domain.  The two
meshes share no grid points; information is exchanged only through
\emph{trilinear interpolation} at selected interface cells.

\begin{notebox}[Recursive nesting]
  \texttt{UpdateDomainInterpolation} is recursive: a fine mesh can
  itself serve as the background for a still-finer mesh, enabling
  cascaded refinement.  The discussion below focuses on the two-level
  case; the recursive extension is straightforward.
\end{notebox}

% ----------------------------------------------------------------
\section{Cell Classification via \texttt{meshTag}}
\label{sec:overset-meshtag}
% ----------------------------------------------------------------

Each cell in every domain carries a scalar field \texttt{meshTag}
(stored in \texttt{mesh->lmeshTag} and \texttt{mesh->meshTag}).
The field is set once during initialisation and remains fixed for
non-moving domains.  Three cell classes are defined:

\begin{center}
\begin{tabular}{llll}
\toprule
\textbf{Classification} & \textbf{Tag range} & \textbf{Predicate} & \textbf{Meaning} \\
\midrule
Calculated cell & $\texttt{meshTag} < 0.1$       & \texttt{isCalculatedCell} & Normal fluid; fully solved \\
Interpolated cell & $0.1 < \texttt{meshTag} < 1.5$ & \texttt{isInterpolatedCell} & Interface layer; value set by overset interpolation \\
Zeroed cell       & $\texttt{meshTag} > 1.5$       & \texttt{isZeroedCell}        & Solid-equivalent blanked region; all fluxes zeroed \\
\midrule
Overset cell      & $\texttt{meshTag} > 0.1$       & \texttt{isOversetCell}       & Union of interpolated + zeroed \\
\bottomrule
\end{tabular}
\end{center}

\paragraph{Background mesh (parent).}
Interior and boundary cells that lie outside the fine-mesh footprint are
\emph{calculated}.  Cells inside the fine-mesh footprint are \emph{zeroed}
(hole-cut blanked): all residuals and fluxes are suppressed there.
A thin collar of cells at the edge of the hole is \emph{interpolated}:
these receive values from the fine mesh (C2P direction) and act as
Dirichlet-like boundary conditions for the background solver.

\paragraph{Fine mesh (child).}
Interior cells are fully \emph{calculated}.
The outermost layer of interior cells (the ghost fringe) is
\emph{interpolated}: values are received from the background mesh (P2C
direction) and override the solver solution, coupling the two domains.

\paragraph{Face classifiers.}
Companion predicates (\texttt{isInterpolatedIFace}, \texttt{isZeroedKFace},
etc.) check the sum of \texttt{meshTag} at the two cells sharing a face;
they are used in \texttt{setBackgroundBC} to set contravariant fluxes at
interface faces.

% ----------------------------------------------------------------
\section{Initialisation (\texttt{InitializeOverset})}
\label{sec:overset-init}
% ----------------------------------------------------------------

Overset initialisation, called once before the time loop, performs five
steps:

\begin{enumerate}
  \item \textbf{Hole object definition.}
        The bounding box (or STL surface) of the fine
        mesh is converted to an IBM-like hole object inside the
        background domain.

  \item \textbf{Tag propagation.}
        The \texttt{meshTag} field is assigned using the hole object:
        interior cells within the hole are tagged as zeroed ($= 3$),
        the adjacent fringe as interpolated ($= 1$).
        Nearby IBM solid cells receive an elevated tag ($= 4 \to 2$ or $3$)
        to prevent incorrect solver behaviour at IBM--overset intersections.

  \item \textbf{P2C acceptor list (\texttt{createAcceptorCellOverset}).}
        For each interpolated cell on the fine mesh, the routine records
        the receiving rank and cell indices, and launches an octree
        donor search inside the background mesh to find the
        enclosing cell.  The result is a list of
        \texttt{aCell} structs (one per acceptor), each storing:
        rank of donor, trilinear weights, and the 8 surrounding
        donor-cell indices.

  \item \textbf{C2P acceptor list (\texttt{createAcceptorCellBackground}).}
        For each interpolated cell on the background mesh, the routine
        identifies the eight corresponding fine-mesh vertices used for the
        C2P averaging stencil (one donor per vertex corner of the
        background cell), again via octree search.

  \item \textbf{Barrier synchronisation.}
        A global \texttt{MPI\_Barrier} after initialisation ensures all
        processes have completed the search before the time loop starts.
\end{enumerate}

% ----------------------------------------------------------------
\section{Per-Step Interpolation (\texttt{UpdateOversetInterpolation})}
\label{sec:overset-update}
% ----------------------------------------------------------------

At the end of each time step, after all domain solvers have advanced to
$t^{n+1}$, the function

\begin{lstlisting}[language=C, style=cstyle]
PetscErrorCode UpdateOversetInterpolation(domain_ *domain);
\end{lstlisting}

is called once from \texttt{main.c}.  It loops over all top-level
domains and calls the recursive helper \texttt{UpdateDomainInterpolation},
which performs interpolation in two passes:

\begin{enumerate}
  \item \textbf{P2C pass (\texttt{interpolateACellTrilinearP2C}).}
        For every interpolated cell on the \emph{fine} mesh, trilinear
        interpolation from the background mesh donor provides updated
        values of $\mathbf{u}$, $T$, and $p$.

  \item \textbf{C2P pass (\texttt{interpolateACellTrilinearC2P}).}
        For every interpolated cell on the \emph{background} mesh, values
        of $\mathbf{u}$, $T$, and $p$ are collected from the eight
        surrounding fine-mesh cells and averaged (spatial average over
        $2\times2\times2$ stencil) to provide coarse-cell values.

  \item \textbf{Scatter.}
        After all MPI communication completes, updated global vectors are
        scattered to their ghost-padded local counterparts
        (\texttt{DMGlobalToLocal}).

  \item \textbf{Recursive descent.}
        If the fine mesh itself hosts a finer child, the procedure is
        applied recursively before returning.
\end{enumerate}

\begin{notebox}[Why 8-vertex averaging for C2P?]
  The C2P direction transfers fine-mesh data to a coarser background
  cell.  Using a single interpolated point (as in P2C) would discard
  spatial variation resolved by the fine mesh.  Averaging over the
  $2\times2\times2$ vertices provides a physically meaningful, low-pass
  filtered boundary condition for the background solver.  It is
  equivalent to a cell-volume weighted injection in the context of
  multigrid prolongation/restriction.
\end{notebox}

% ----------------------------------------------------------------
\section{Variable Linkage}
\label{sec:overset-variables}
% ----------------------------------------------------------------

\subsection{Transferred fields}

\begin{center}
\begin{tabular}{lllll}
\toprule
\textbf{Variable} & \textbf{Direction} & \textbf{Donor source} & \textbf{Acceptor target} & \textbf{MPI tag} \\
\midrule
\texttt{Ucat} (Cartesian $\mathbf{u}$) & P2C & \texttt{ueqnD->lUcat}  & \texttt{ueqnA->Ucat} $\to$ \texttt{lUcat} & 0 \\
\texttt{Tmprt} (pot.\ temperature $T$) & P2C & \texttt{teqnD->lTmprt} & \texttt{teqnA->Tmprt} $\to$ \texttt{lTmprt} & 1 \\
\texttt{P} (pressure $p$)              & P2C & \texttt{peqnD->lP}     & \texttt{peqnA->P} $\to$ \texttt{lP}       & 2 \\
\midrule
\texttt{Ucat} (8-avg) & C2P & \texttt{ueqnD->lUcat}  & \texttt{ueqnA->Ucat} $\to$ \texttt{lUcat} & $b$ \\
\texttt{Tmprt} (8-avg) & C2P & \texttt{teqnD->lTmprt} & \texttt{teqnA->Tmprt} $\to$ \texttt{lTmprt} & $b{+}N_A$ \\
\texttt{P} (8-avg)    & C2P & \texttt{peqnD->lP}     & \texttt{peqnA->P} $\to$ \texttt{lP}       & $b{+}2N_A$ \\
\bottomrule
\end{tabular}
\end{center}

Here $b$ is a base tag offset that differs per domain pair, and $N_A$ is the
number of acceptor cells on that rank.
The three variables use consecutive tag offsets within each pass,
ensuring no tag collision between simultaneous \texttt{MPI\_Send}/\texttt{MPI\_Recv}
calls.

\subsection{How the transferred fields feed into each solver}

\paragraph{Velocity — \texttt{Ucat} / \texttt{lUcat}.}
After the scatter, \texttt{setBackgroundBC} is called at the end of both
P2C and C2P interpolation functions.
For every interpolated \emph{face}, it projects the ghost-layer Cartesian
velocity stored in \texttt{lUcat} onto the local face-area vectors to obtain
contravariant fluxes \texttt{ucont[k][j][i].x/y/z}:

\begin{lstlisting}[style=cstyle, language=C, caption={\texttt{setBackgroundBC} projects interpolated Cartesian velocity to contravariant fluxes.}]
if (isInterpolatedIFace(k, j, i, i+1, meshTag))
{
    ucont[k][j][i].x =
        0.5*(  ucat[k][j][i  ].x * icsi[k][j][i  ].x
             + ucat[k][j][i  ].y * icsi[k][j][i  ].y
             + ucat[k][j][i  ].z * icsi[k][j][i  ].z
             + ucat[k][j][i+1].x * icsi[k][j][i+1].x
             + ... );
}
\end{lstlisting}

This is the primary mechanism by which the overset boundary condition enters
\texttt{FormU}: the contravariant flux at an interface face carries the
interpolated velocity from the neighbouring domain, not the locally computed
value.  For zeroed cells, the flux is set to zero (\texttt{isZeroedIFace}
branch), suppressing all convective and viscous contributions across the
hole boundary.

\paragraph{Pressure — \texttt{P} / \texttt{lP}.}
The pressure field at interface and blanked cells is fully controlled by
the overset interpolation; the Poisson solver does \emph{not} modify it.
This is guaranteed by the trivial-row construction in
\texttt{SetCoeffMatrix}: every overset cell (\texttt{isOversetCell})
receives the identity equation $A_{ii}=1$, $A_{ij}=0$ $(j\neq i)$, so
the Poisson solver sets $\phi=0$ for those cells.
After \texttt{UpdatePressure} adds $\phi=0$ to $p$, the pressure at
interface cells is exactly the value deposited by the previous overset
interpolation step.

\texttt{GradP} reads \texttt{lP} across the interface (including
ghost cells at interpolated and zeroed positions) to compute the
pressure-gradient term $\nabla p$ for the momentum equation.
Because \texttt{lP} is fully populated by \texttt{UpdateOversetInterpolation}
before the next step's \texttt{SolveUEqn} is called, \texttt{GradP}
receives consistent, physically meaningful values at every stencil position.

\paragraph{Temperature — \texttt{Tmprt} / \texttt{lTmprt}.}
The temperature at interface cells is deposited by overset interpolation.
\texttt{FormT} evaluates convective and diffusive fluxes across the interface
using the ghost-layer values in \texttt{lTmprt}, thereby coupling thermal
transport between the two domains.
Blanked cells receive $T=0$ from the zeroed-flux treatment, ensuring
no spurious heat source.

\begin{warnbox}[One-step lag]
  \texttt{UpdateOversetInterpolation} is called \emph{after} all domain
  solvers have advanced to $t^{n+1}$.  The interpolated values deposited
  into \texttt{lUcat}, \texttt{lP}, and \texttt{lTmprt} therefore correspond
  to state $t^{n+1}$ of the \emph{donor} domain, and are consumed by the
  \emph{acceptor} domain on the \emph{next} time step ($t^{n+1}$ to
  $t^{n+2}$).  This introduces a one-step lag between coupled domains.
  For well-resolved, low-CFL simulations the lag is negligible; for large
  CFL or tightly coupled problems it may require implicit coupling
  (not currently implemented).
\end{warnbox}

% ----------------------------------------------------------------
\section{Position in the Time Loop}
\label{sec:overset-timeloop}
% ----------------------------------------------------------------

The following pseudocode shows the call order inside the main time loop
for each domain \texttt{d}:

\begin{lstlisting}[style=cstyle, language=C, caption={Simplified main time loop with overset calls (src/main.c).}]
for (each time step)
{
    for (d = 0; d < nDomains; d++)
    {
        Buoyancy(ueqn, T^n);          // buoyancy from lTmprt (prev step)

        SolveUEqn(ueqn);              // uses lUcat, lP from prev overset step
        SolvePEqn(peqn);              // phi=0 at overset cells; lP unchanged there
        contravariantToCartesian(ueqn);

        SolveTEqn(teqn);              // uses lTmprt from prev overset step

        UpdateCartesianBCs(ueqn);
        UpdateContravariantBCs(ueqn);
    }

    if (isOversetActive)
        UpdateOversetInterpolation(domain); // deposit lUcat, lP, lTmprt for next step
}
\end{lstlisting}

Key observations:
\begin{itemize}
  \item \texttt{SolveUEqn} and \texttt{SolvePEqn} consume the
        overset-interpolated values from the \emph{previous} call to
        \texttt{UpdateOversetInterpolation}.

  \item \texttt{UpdateOversetInterpolation} is called once per global step,
        \emph{after} all domains have been advanced, not interleaved.
        The recursive descent inside \texttt{UpdateDomainInterpolation}
        ensures parent-to-child (P2C) is applied before
        child-to-parent (C2P) within each level.

  \item The Poisson solver correctly ignores overset cells: their identity
        rows force $\phi=0$, so \texttt{UpdatePressure} leaves their
        pressure unchanged, and the only update to $p$ at those cells
        is the one performed by \texttt{UpdateOversetInterpolation}.
\end{itemize}

% ----------------------------------------------------------------
\section{Known Limitations}
\label{sec:overset-limitations}
% ----------------------------------------------------------------

\begin{itemize}
  \item \textbf{Load-set interpolation not implemented.}
        Only trilinear (8-point) interpolation is available.
        Higher-order or least-squares schemes are not currently supported.

  \item \textbf{Rotating meshes not implemented.}
        The acceptor lists are built assuming fixed mesh positions.
        Mesh rotation or translation would require rebuilding the
        acceptor lists every step (expensive) or using a moving-mesh
        framework not currently present.

  \item \textbf{One-step coupling lag.}
        See the warning box in Section~\ref{sec:overset-variables}.
        Tightly coupled phenomena (e.g.\ rapidly oscillating pressure
        waves near the interface) are limited by this lag.

  \item \textbf{Poisson null-space at interface.}
        Because interface cells contribute $\phi=0$ (identity rows),
        they are excluded from the divergence-free projection.
        Small divergence errors may accumulate at the overset fringe
        over long runs; it is recommended to monitor continuity errors
        at overset boundaries.
\end{itemize}

\end{document}


